% =====================================================================
%  capitulos/cap13.tex — Capítulo 13. Las relaciones interoracionales:
%  subordinación
%  Autores: Augusto M. Trombetta y Mabel Giammatteo
%  Fuente: Capitulo_13_normalizado.docx
%
%  Siglas (decisión editorial, ver HANDOFF): Ocomp = oración compuesta;
%  OS = oración subordinada (genérica); OC = oración completiva;
%  OR = oración relativa.
% =====================================================================
\capitulo{Las relaciones interoracionales: subordinación}{Augusto M. Trombetta y Mabel Giammatteo}

\section{Introducción}

En este capítulo se completa el estudio de las relaciones interoracionales con el
tratamiento de la subordinación y sus principales subclases. En primer lugar, se
plantea en qué consiste la subordinación oracional (la oración carente de
autonomía sintáctica, esto es, incluida en una estructura
mayor)\pendiente{El original dice «no incluida en una estructura mayor»} y,
luego, se hace foco en la estructura interna de las oraciones subordinadas. Para
ello se tiene en cuenta la presencia (o no) de subordinantes léxicos y su
naturaleza o bien conjuntiva o bien relativa, así como el tipo de flexión verbal
(conjugada o no finita) involucrada. La consideración de estos aspectos
estructurales se completa con el tratamiento de las funciones sintácticas que
cumplen las oraciones subordinadas. La tradición gramatical clasificó a las
subordinadas teniendo en cuenta su equivalencia funcional con algunas clases de
palabras, por lo que se distinguía entre subordinadas sustantivas, adjetivas y
adverbiales. Sin embargo, aquí se adopta una clasificación que, según el tipo de
encabezador que las introduce, distingue entre subordinadas completivas y
relativas y, dentro de estas últimas, diferencia las que se ligan con un elemento
nominal o adverbial previo, que funciona como su antecedente, y las que carecen
de antecedente léxico y se conocen como relativas libres y semilibres.
Finalmente, desde la perspectiva semántico-pragmática se consideran las
construcciones que incluyen subordinadas que transmiten valores
lógico-relacionales (causales, finales, concesivas y condicionales) o bien
valores cuantificacionales-intensivos (comparativas y consecutivas).

\section{La relación de subordinación}\label{sec:13-relacion}

Según plantean \textcite[731]{bosque2009}, «[l]a subordinación es una relación de
dependencia sintáctica», en la que se establece una relación jerárquica entre dos
constituyentes sintácticos, uno de los cuales complementa o modifica al otro, que
funciona como núcleo. Así entendida, la relación de subordinación se aplica por
igual a cualquier vinculación entre un núcleo y sus complementos o adjuntos, con
independencia de la configuración sintáctica que adopten:

\begin{cita}
la relación sintáctica que existe entre el verbo y su \textit{complemento} es la
misma en \textit{Lamento que las cosas estén así} que en \textit{Lamento la
situación}, aun cuando tradicionalmente se entendía que la oración que encabeza
\textit{que} estaba \sig{subordinada}, mientras que esa noción no se aplicaba al
grupo nominal \textit{la situación} \parencite[§\,1.13m]{rae2025}.
\end{cita}

\begin{destacado}
\textbf{Para tener en cuenta.} En sintonía con lo anterior, se puede señalar que
la ampliación del concepto de \termino{nuclearidad} a todas las categorías
gramaticales (incluidas las preposiciones; véase el \capref{5},
§\ref{sec:05-sp}), junto con la universalización de la
\termino{endocentricidad} (que sostiene que toda construcción gramatical posee un
núcleo; véase el \capref{5}, §\ref{sec:05-sintagmas}), ha permitido entender que
el núcleo de una construcción gramatical es el constituyente al que se
subordinan los restantes constituyentes.
\end{destacado}

\subsection{La subordinación oracional}\label{sec:13-oracional}

Centrándonos ahora en las relaciones entre oraciones, la
\termino{subordinación} (o \termino{hipotaxis}, según la denominación de las
gramáticas clásicas) es la relación estructural entre una oración subordinada y
una oración matriz (también llamada \textit{oración principal}), a la que se
integra y de la que forma parte. En consecuencia, la subordinada queda incluida
en un constituyente más extenso, la oración principal
\parencite[214--215]{bosque2019}, y forma con ella una estructura mayor conocida
como \termino{oración compuesta}.\footnote{El concepto de oración compuesta se
presentó en el \capref{12}, §\ref{sec:12-pasaje}.} Según iremos viendo, la
subordinada «desempeña una función dentro de la oración "principal" (sujeto,
complemento directo, etc.) o dentro de un constituyente de ella»
\parencite[731]{bosque2009}. El esquema de \exref{ex:13-diag1} representa la
relación.

\begin{exe}
  \ex\label{ex:13-diag1}
      \begin{forest} gramatica
        [Ocomp\\(oración matriz)
          [SN\\Suj]
          [SV\\Pred
            [V]
            [OS\\(argumento o adjunto) [\ldots, roof]]]]
      \end{forest}
\end{exe}

La oración subordinada (OS) se encuentra inserta en una estructura mayor
(oración matriz) de la que forma parte; junto con ella, forma una oración
compuesta (Ocomp). Dentro de esa estructura, la subordinada establece una
relación de dependencia y puede cumplir una función argumental (sujeto,
complemento) o no argumental (adjunto). En corchetes: [\etq{Ocomp} \ldots\
[\etq{OS} \ldots] ].\pendiente{Diagrama 1 rehecho («cambiar diagrama»)}

La subordinación involucra una diferencia sintáctica entre la oración subordinada
y la oración matriz a la que se integra, ya que la subordinada es una estructura
carente de autonomía sintáctica. En cambio, la oración compuesta, que resulta de
la relación matriz-subordinada, constituye una estructura sintácticamente
autónoma:

\begin{exe}
  \ex\label{ex:13-gobierno}
    \begin{xlist}
      \ex\label{ex:13-gobierno-a} {}[\etq{Ocomp} El Gobierno decidió [\etq{OC}
          aumentar los sueldos de las fuerzas de seguridad]].
      \ex\label{ex:13-gobierno-b} {}[\etq{Ocomp} El Gobierno decidió [\etq{OC}
          que aumenten los sueldos de las fuerzas de seguridad]].
    \end{xlist}
\end{exe}

Tanto \exref{ex:13-gobierno-a} como \exref{ex:13-gobierno-b} son oraciones
compuestas, que constituyen estructuras completas en términos
sintáctico-semánticos y pueden ser inicio de discurso. Ambas cualidades son
aplicables a cualquier oración simple por la que se puedan intercambiar:

\begin{exe}
  \ex\label{ex:13-simples}
    \begin{xlist}
      \ex El Gobierno decidió [\etq{SN} el aumento de los sueldos de las fuerzas
          de seguridad].
      \ex El Gobierno decidió [\etq{SN} un aumento de sueldo para las fuerzas de
          seguridad].
    \end{xlist}
\end{exe}

En contraste, aun cuando constituye una estructura completa en términos
sintáctico-semánticos, puesto que contiene una predicación, la subordinada no
puede ser inicio de discurso:

\begin{exe}
  \ex
    \begin{xlist}
      \ex[??]{aumentar los sueldos de las fuerzas de seguridad.}
      \ex[??]{que aumenten los sueldos de las fuerzas de seguridad.}
    \end{xlist}
\end{exe}

En cuanto a la oración matriz, puede ser una estructura completa, como las de
\exref{ex:13-simples}, o la presencia de la subordinada puede ser obligatoria
porque cumple una función argumental (sujeto, complemento) en ella. En cambio,
cuando la subordinada no es argumental, funciona como un modificador adjunto
dentro de la oración matriz:

\begin{exe}
  \ex\label{ex:13-decidio}
    \begin{xlist}
      \ex\label{ex:13-decidio-a} El Gobierno decidió [\etq{OC} que aumenten los
          sueldos de las fuerzas armadas].
      \ex\label{ex:13-decidio-b} El Gobierno decidió [\etq{OC} aumentar los
          sueldos de las fuerzas armadas].
      \ex\label{ex:13-decidio-c} El Gobierno decidió algo inesperado, [\etq{OC}
          que aumenten los sueldos de las fuerzas armadas].
      \ex\label{ex:13-decidio-d} El Gobierno decidió algo inesperado, [\etq{OC}
          aumentar los sueldos de las fuerzas armadas].
    \end{xlist}
\end{exe}

En \exref{ex:13-decidio-a} y \exref{ex:13-decidio-b}, la oración subordinada
cumple la función argumental de complemento objeto directo del verbo transitivo
\forma{decidir} ⟨alguien, algo⟩. En ese sentido, la subordinada contribuye a la
buena formación (gramaticalidad; véase el \capref{2},
§\ref{sec:02-aceptabilidad}) de la oración matriz, ya que
\juicio{*}{\forma{El Gobierno decidió}} forma una secuencia incompleta. En
\exref{ex:13-decidio-c} y \exref{ex:13-decidio-d}, en cambio, la misma oración
subordinada cumple la función no argumental de aposición explicativa de un
núcleo nominal (el pronombre neutro \forma{algo}; véase el \capref{5},
§\ref{sec:05-sn}) y su eventual ausencia no afectaría la buena formación
oracional, ya que \forma{El Gobierno decidió algo inesperado} es una oración
completa, tanto desde el punto de vista sintáctico como semántico.

Asimismo, en un sentido formal, los ejemplos de \exref{ex:13-decidio} dejan ver
que, internamente, las oraciones subordinadas pueden constituirse alrededor de un
verbo conjugado en forma personal (\forma{El Gobierno decidió [que
\textbf{aumenten} los sueldos de las fuerzas armadas]}) o de un verbo en forma no
personal (\forma{El Gobierno decidió [\textbf{aumentar} los sueldos de las
fuerzas armadas]}). En forma correlativa, las subordinadas con verbos finitos
están encabezadas por un subordinante que puede ser conjuntivo o relativo,
mientras que las subordinadas con verbos no finitos carecen de un encabezador que
destaque la relación subordinativa.

Recordemos que en el \capref{6}, §\ref{sec:06-capas}, se introdujo la idea de que
la estructura de la oración se puede entender como un conjunto de tres capas
formales (complementante (SC), flexiva (SF) y léxica (SV)) que se corresponden
con valores discursivos, deícticos y semánticos, respectivamente. El esquema de
\exref{ex:13-diag2} muestra que el núcleo complementante (C) puede estar
realizado léxicamente por una conjunción \forma{que / si} o puede no tener
realización léxica (Ø)\footnote{El símbolo Ø indica una categoría
fonológicamente vacía, es decir, que no está realizada por un elemento léxico.}
en la oración subordinada.

\begin{exe}
  \ex\label{ex:13-diag2}
      \begin{forest} gramatica
        [{SC = OS}
          [C\\\textit{que / si / Ø}]
          [SF [\ldots, roof]]]
      \end{forest}
\end{exe}

Dentro de la oración subordinada, los subordinantes conjuntivos (expresos o
nulos) encabezan el sintagma complementante. Como corresponde a su función
nuclear, el subordinante determina la naturaleza de la oración subordinada
(conjuntiva; véase §\ref{sec:13-encabezador}). El sintagma complementante es la
máxima proyección oracional (es decir, SC = O) y su complemento es el SF.

Así, es importante destacar que la relación de subordinación puede estar presente
aun en ausencia de encabezadores que hagan explícito el vínculo entre principal y
subordinada:

\begin{exe}
  \ex\label{ex:13-xi}
    \begin{xlist}
      \ex\label{ex:13-xi-a} Xi Jinping pregunta [\etq{OC} cuánta presión de
          Estados Unidos soportará China].
      \ex\label{ex:13-xi-b} Xi Jinping busca [\etq{OC} ejercer presión sobre
          Estados Unidos].
    \end{xlist}
\end{exe}

En los ejemplos de \exref{ex:13-xi}, ni la subordinada interrogativa
\exref{ex:13-xi-a} ni la subordinada de infinitivo \exref{ex:13-xi-b} poseen un
introductor subordinante; sin embargo, la relación de subordinación se sostiene
por los requerimientos del verbo selector correspondiente. Ambos ejemplos, las
subordinadas de naturaleza enfática (exclamativas e interrogativas parciales;
véase la nota~\ref{fn:13-exclamativas}) y las subordinadas con verbo no finito
(infinitivo, gerundio y participio), son casos típicos de oraciones subordinadas
con subordinante nulo (o sin expresión léxica). En el caso de las subordinadas
enfáticas, los ítems exclamativos o interrogativos iniciales (\forma{cuánta,
dónde, cómo}, etc.) no se ubican en el núcleo complementante, sino que
desempeñan su función sintáctica dentro de la oración subordinada de la que
forman parte. Por ejemplo, \forma{cuánta} es el especificador del SN \forma{presión
de Estados Unidos}, es decir que cumple la misma función sintáctica que le
correspondería en una oración exclamativa o interrogativa no subordinada (por
ejemplo, \forma{¿Cuánta presión de Estados Unidos soportará China?}), como se ve
en \exref{ex:13-diag3}.

\begin{exe}
  \ex\label{ex:13-diag3}
      \begin{forest} gramatica
        [{SC = OC}
          [SN\\Foco
            [D\\\textit{cuánta}]
            [N\prima [\textit{presión de Estados Unidos}, roof]]]
          [C\prima
            [C\\Ø]
            [SF [\textit{soportará China}, roof]]]]
      \end{forest}
\end{exe}

Por su lado, las oraciones subordinadas de formas verbales no finitas (verbos en
formas no personales) carecen de encabezadores léxicos que determinen la
naturaleza de la relación subordinativa \exref{ex:13-diag4}.

\begin{exe}
  \ex\label{ex:13-diag4}
      \begin{forest} gramatica
        [{SC = OC}
          [C\\Ø]
          [SF
            [SN\\PRO]
            [F\prima
              [F\\\textit{-r} (inf.)]
              [SV [\textit{ejercer presión sobre Estados Unidos}, roof]]]]]
      \end{forest}
\end{exe}

\pendiente{Diagramas 3 y 4: el original no conserva la imagen; se reconstruyeron}

\begin{ejercicio}\label{ej:13-paremias}
En muchas paremias (refranes, dichos y proverbios tradicionales) suelen aparecer
oraciones subordinadas. Identifique en los siguientes ejemplos la oración
subordinada (OS) dentro de la oración compuesta (Ocomp) y colóquela entre
corchetes, según se muestra en los ejemplos resueltos. Tenga en cuenta que puede
existir más de una oración subordinada en una matriz.

\begin{exe}
  \ex A la cama no te irás sin saber una cosa más.
  \ex A quien madruga, Dios lo ayuda.
  \ex Cuando el invierno primavea, la primavera invernea.
  \ex Cuando veas las barbas de tu vecino cortar, pon las tuyas a remojar.
  \ex Dando gracias por agravios negocian los hombres sabios.
  \ex Diciendo las verdades se pierden amistades.
  \ex Donde fueres, haz lo que vieres.
  \ex Donde hubo fuego, cenizas quedan.
  \ex El salir de la posada es la mejor jornada.
  \ex Es de bien nacidos ser agradecidos.
  \ex Hablando del rey de Roma, por la puerta asoma.
  \ex Lo que natura non da Salamanca non presta.
  \ex No por mucho madrugar amanece más temprano.
  \ex Nunca es tarde si la dicha es buena.
  \ex Pájaro que comió, voló.
  \ex Perro que ladra no muerde.
  \ex Preguntando se llega a Roma.
  \ex Si quieres ser bien servido, sírvete a ti mismo.
\end{exe}

\begin{resuelto}[Ejemplos resueltos]
1. [\etq{Ocomp} A la cama no te irás sin [\etq{OS} \textit{saber una cosa más}]].
2. [\etq{Ocomp} A [\etq{OS} \textit{quien madruga}], Dios lo ayuda].
\end{resuelto}
\end{ejercicio}

\section{Estructura interna de las oraciones subordinadas}\label{sec:13-estructura}

Dentro de este apartado atendemos a dos cuestiones estructurales: en primer
lugar, damos cuenta de la presencia o no de un introductor de la subordinada y
del tipo de encabezador que la introduce ---conjuntivo o relativo---; en segundo
lugar, nos ocupamos del tipo de flexión ---conjugada (finita o personal) o no
conjugada (no finita o no personal)--- que manifiesta el verbo de la subordinada.

\subsection{Presencia y tipo de encabezador}\label{sec:13-encabezador}

Las oraciones subordinadas de verbo finito (o conjugado en formas personales) son
introducidas por un ítem léxico, una conjunción subordinante (\forma{que, cuando,
si}, etcétera). Recordando lo visto en el \capref{5}, §\ref{sec:05-clases}, sobre
las clases de palabras funcionales o gramaticales, las conjunciones son unidades
léxicas que establecen relaciones entre distintos constituyentes de la oración
para conectarlos sintácticamente \parencite{giammatteo2009}. Esta función
gramatical de conexión involucra dos tipos diferentes de relaciones, según la
jerarquía sintáctica de los constituyentes involucrados: la coordinación, tratada
en el capítulo anterior, y la subordinación. Se reconocen así dos grandes
subclases de ítems conjuntivos: los coordinantes y los subordinantes. Las
conjunciones coordinantes, según lo ya visto, establecen relaciones entre
constituyentes de igual jerarquía sintáctica. Los elementos conjuntivos de
naturaleza subordinante, en contraste, marcan relaciones entre constituyentes de
distinta jerarquía sintáctica y tienen como función incluir una estructura
predicativa, la subordinada, en otra, que es la oración principal o matriz.

Asimismo, dentro de los subordinantes, se pueden reconocer dos tipos o subclases:
las conjunciones y los relativos. Si las conjunciones subordinantes (tales como
\forma{que, cuando, si}) se pueden caracterizar a partir de su función
subordinativa, los relativos (tales como \forma{que, el que, quien, cual, cuyo,
donde, cuando, como, cuanto}), que son de naturaleza pronominal, además de
incluir una oración subordinada en la matriz, cumplen una función sintáctica
determinada dentro de ella \parencite{giammatteo2009}.

Las conjunciones subordinantes, además de hacer explícita la relación de
subordinación, pueden aportar valores relacionados con la modalidad de la oración
subordinada. En los ejemplos de \exref{ex:13-modalidad}, las conjunciones
subordinantes marcan la modalidad declarativa \exref{ex:13-modalidad-a},
interrogativa \exref{ex:13-modalidad-b} o exhortativa \exref{ex:13-modalidad-c}
de las oraciones subordinadas que encabezan, en consonancia con los
requerimientos del verbo matriz correspondiente:

\begin{exe}
  \ex\label{ex:13-modalidad}
    \begin{xlist}
      \ex\label{ex:13-modalidad-a} Xi Jinping afirma [\etq{OC} \textit{que} la
          presión de Estados Unidos sobre China no tendrá éxito].
      \ex\label{ex:13-modalidad-b} Xi Jinping pregunta [\etq{OC} \textit{si} la
          presión de Estados Unidos sobre China tendrá éxito].
      \ex\label{ex:13-modalidad-c} Xi Jinping ordenó [\etq{OC} \textit{que} sus
          funcionarios no cedan ante la presión de Estados Unidos].
    \end{xlist}
\end{exe}

Los relativos, al igual que las conjunciones subordinantes, hacen explícita la
relación de subordinación y, además, suman una segunda función, interna a la
subordinada que introducen. En los ejemplos de \exref{ex:13-relativos}, los
relativos \forma{que} \exref{ex:13-relativos-a} y \forma{donde}
\exref{ex:13-relativos-b} cumplen funciones sintácticas en la oración
subordinada. En \exref{ex:13-relativos-a}, el relativo, cuyo antecedente es
\forma{la presión}, es el complemento objeto directo. Por su parte, en
\exref{ex:13-relativos-b}, el relativo, cuyo antecedente es \forma{las áreas}, es
el complemento de la preposición \forma{de}, con la que forma un complemento
circunstancial de lugar:

\begin{exe}
  \ex\label{ex:13-relativos}
    \begin{xlist}
      \ex\label{ex:13-relativos-a} Xi Jinping condena la presión [\etq{OR}
          [\textit{que}]\etq{COD} ejerce Estados Unidos sobre China].\\
          \textit{que} = la \textit{presión}.
      \ex\label{ex:13-relativos-b} Xi Jinping aumenta las áreas [\etq{OR} [de
          \textit{donde}]\etq{CCL} China extrae hierro y carbón].\\
          (de) \textit{donde} = (de) las \textit{áreas}.
    \end{xlist}
\end{exe}

Como ya adelantamos, típicamente carecen de introductor léxico explícito las
oraciones subordinadas articuladas alrededor de ítems interrogativos y
exclamativos, y también otras subordinadas que admiten la inserción de la
conjunción subordinante \forma{que}, así como las subordinadas articuladas
alrededor de un verbo no finito (o no conjugado en formas personales). En
síntesis, las oraciones subordinadas se pueden articular de dos modos: con o sin
introductor léxico (tabla~\ref{tab:13-encabezadores}).

{\small
\begin{longtable}{@{}>{\raggedright\arraybackslash}p{40mm}
    >{\raggedright\arraybackslash}p{83mm}@{}}
  \caption{Subordinación con y sin encabezadores léxicos}\label{tab:13-encabezadores}\\
  \toprule
  \endfirsthead
  \toprule
  \endhead
  Subordinación con encabezadores léxicos explícitos (marcan relaciones entre
    constituyentes de distinta jerarquía sintáctica) &
    \textbf{Conjunciones subordinantes.} Semánticamente, indican modalidad (son
    índices de modalidad declarativa, interrogativa, exhortativa).
    Sintácticamente, encabezan el sintagma complementante (SC), sin cumplir otra
    función sintáctica adicional. Morfológicamente, son invariables. \\
  \addlinespace
  & \textbf{Relativos.} Semánticamente, toman la referencia de un antecedente o
    suman valores generales (cosa, persona, cantidad, lugar, modo, tiempo).
    Sintácticamente, introducen el sintagma complementante (SC).
    Morfológicamente, pueden ser invariables o variables. \\
  \midrule
  Subordinación sin encabezadores léxicos explícitos &
    Oraciones subordinadas interrogativas y exclamativas. Oraciones subordinadas
    con \textit{que} facultativo. Oraciones subordinadas con formas verbales no
    finitas. \\
  \bottomrule
\end{longtable}}

\begin{ejercicio}\label{ej:13-introductor}
Retome los resultados del ejercicio~\ref{ej:13-paremias} y, si corresponde,
señale el tipo de introductor de cada subordinada, según se muestra en los
ejemplos resueltos.

\begin{resuelto}[Ejemplos resueltos]
1. Sin subordinante. 2. Con relativo \forma{quien}.
\end{resuelto}
\end{ejercicio}

\begin{ejercicio}\label{ej:13-principito}
A partir de los siguientes pares de oraciones simples, adaptadas de \textit{El
principito}, de Antoine de Saint-Exupéry (1943), construya una oración compuesta
(matriz-subordinada), según se muestra en el ejemplo resuelto~1.

\begin{exe}
  \ex Lo esencial es invisible a los ojos. Lo dijo el Principito.
  \ex A veces tenemos que aguantar a las orugas. Queremos disfrutar de las
      mariposas.
  \ex Eres el dueño de tu vida y tus emociones. Nunca te olvides de esto.
  \ex Los hombres ya no tienen tiempo. Compran cosas hechas a los comerciantes.
  \ex Me pregunto algo. ¿Las estrellas se iluminan con algún fin?
  \ex Nunca estamos contentos en un lugar. Estamos en ese lugar.
  \ex Todas las estrellas ríen. En ese momento brilla la luna.
  \ex Todas las personas grandes han sido niños antes. Pocas lo recuerdan.
  \ex Un hombre es vanidoso. Todos los demás hombres son admiradores.
  \ex Uno termina de arreglarse. Entonces hay que hacer cuidadosamente la
      limpieza del planeta.
\end{exe}

\begin{resuelto}
1. [\etq{Ocomp} El Principito dijo [\etq{OC} que lo esencial era invisible a los
ojos]].
\end{resuelto}
\end{ejercicio}

\subsubsection{Los encabezadores de las subordinadas}\label{sec:13-encabezadores}

Según hemos dicho, hay dos tipos de encabezadores léxicos de oraciones
subordinadas: las conjunciones subordinantes y los relativos. A continuación,
dedicamos un apartado a cada uno de estos tipos de introductores.

\paragraph{Las conjunciones subordinantes.}\label{sec:13-conjunciones}
Las conjunciones subordinantes son los núcleos funcionales de la oración (es
decir, del sintagma complementante, la máxima proyección oracional) y determinan
la naturaleza de la subordinada. Carentes de valor referencial, son responsables
de la interpretación modal de la subordinada; así, según se mostró en los
ejemplos de \exref{ex:13-modalidad}, \forma{que} es índice de modalidad
declarativa / exhortativa y \forma{si} marca modalidad interrogativa.

En los análisis clásicos de la subordinación oracional, las conjunciones
subordinantes y las locuciones conjuntivas derivadas de estas se asociaban a
distintos valores lógico-relacionales (tiempo, causa, concesión, condición,
finalidad). En otras palabras, se contaba con una lista de formas conjuntivas
simples y complejas dotadas de valores específicos: \forma{porque, como, que, a
causa de que, en razón de que, en la medida en que, ya que, puesto que, dado
que} (causales); \forma{aunque, aun cuando, si bien, eso que, a pesar de que,
pese a que, aun si, incluso si} (concesivas); \forma{si, de, siempre que, siempre
y cuando, a no ser que, en caso de que, a condición de que, con tal de que}
(condicionales); \forma{para que, a fin de que, con el objeto de que, con el fin
de que, con vistas a que} (finales); \forma{antes que, después que, siempre que,
ni bien, no bien} (temporales).

En la enumeración anterior, además de ciertas repeticiones de formas, se advierte
fácilmente que las conjunciones y locuciones subordinantes, en su gran mayoría,
pueden resumirse como una construcción encabezada por un adverbio o preposición
seguidos de una oración subordinada encabezada por las conjunciones \forma{que}
(mayoritariamente) o \forma{si} (en menor medida): \forma{antes (que), después
(que), siempre (que); por(que), a causa de (que), en razón de (que), en la medida
en (que), ya (que), puesto (que), dado (que); aun(que), (si) bien, eso (que), a
pesar de (que), pese a (que), aun (si), incluso (si); siempre (que), a no ser
(que), en caso de (que), a condición de (que), con tal de (que); para (que), a fin
de (que), con el objeto de (que), con el fin de (que), con vistas a (que)}. Por
tanto, si aceptamos la distinción que acabamos de hacer, se pueden admitir dos
versiones en el análisis de las conjunciones y locuciones conjuntivas
subordinantes:

\begin{itemize}
  \item una \termino{versión sintética}, que es la tradicional, en la que las
        locuciones conjuntivas subordinantes funcionan como encabezadores
        unitarios dotados de valores específicos ---\forma{aunque} (concesivo),
        \forma{porque} (causal), \forma{siempre que} (temporal)---;
  \item una \termino{versión analítica}, que es la alternativa propuesta en
        \textcite{bosque2009} y en \textcite{brucart2009,brucart2016}, en la que
        las conjunciones subordinantes funcionan como encabezadores de oraciones
        subordinadas dependientes de un núcleo preposicional o adverbial, que es
        el que posee el valor específico.
\end{itemize}

En forma esquemática, es posible representar la forma \forma{porque} de las
siguientes maneras:

\begin{exe}
  \ex\label{ex:13-porque}
    \begin{xlist}
      \ex\label{ex:13-porque-a} Versión sintética:\\[2pt]
          \begin{forest} gramatica
            [{SC = OS}
              [C\\\textit{porque}]
              [SF [\ldots, roof]]]
          \end{forest}
      \ex\label{ex:13-porque-b} Versión analítica:\\[2pt]
          \begin{forest} gramatica
            [SP
              [P\\\textit{por}]
              [{SC = OC}
                [C\\\textit{que}]
                [SF [\ldots, roof]]]]
          \end{forest}
    \end{xlist}
\end{exe}

En la versión sintética \exref{ex:13-porque-a}, la locución conjuntiva, como
cualquier conjunción subordinante, encabeza el sintagma complementante (SC), la
máxima proyección oracional; es el análisis clásico más extendido: [\etq{O}
\ldots\ [\etq{OS} \textit{porque} \ldots]]. En la versión analítica
\exref{ex:13-porque-b}, la conjunción subordinante encabeza el sintagma
complementante (SC), y esta proyección oracional es el complemento de un núcleo
preposicional o adverbial; es el análisis alternativo más reciente: [\etq{O}
\ldots\ [\etq{SP} \textit{por} [\etq{OC} \textit{que} \ldots]]].

Pese a ser el más extendido, el enfoque sintético de las conjunciones y
locuciones conjuntivas subordinantes pierde de vista una interesante reflexión
respecto de estas estructuras: no parece tomar en cuenta la sistemática aparición
de los dos subordinantes básicos (\forma{que, si}) en todas ellas. Esta
regularidad queda mejor analizada en el enfoque analítico, más simple y
regularizador (aunque inicialmente pueda parecernos más extraño porque no estamos
acostumbrados a pensar así estas construcciones): las conjunciones subordinantes
(\forma{que, si}) encabezan oraciones subordinadas que funcionan como
complementos de adverbios o de preposiciones que, a su vez, son las responsables
del valor que toma la construcción (temporal, causal, concesivo, condicional,
final). Así analizadas, estas expresiones adverbiales y preposicionales son las
mismas que, mayoritariamente, pueden tomar complementos no oracionales para
articular esos mismos valores: \forma{antes (de) \textbf{que comamos} / antes de
\textbf{la comida}; después (de) \textbf{que me visitaste} / después de
\textbf{tu visita}; por\textbf{que me lo pediste} / por \textbf{tu pedido}; a
causa de \textbf{que nevó copiosamente} / a causa de \textbf{una copiosa
nevada}; en razón de \textbf{que intervino el rector} / en razón de \textbf{la
intervención del rector}; dado \textbf{que cumple con los requisitos} / dado
\textbf{el cumplimiento de los requisitos}; aun\textbf{que coma poco} / aun
\textbf{poco alimentado}; a pesar de \textbf{que insistas} / a pesar de
\textbf{tu insistencia}; pese a \textbf{que se marchó} / pese a \textbf{su
salida}; aun \textbf{si estaba enferma} / aun \textbf{enferma}; en caso de
\textbf{que suceda un incendio} / en caso de \textbf{incendio}; para \textbf{que
estudies} / para \textbf{tus estudios}; con vistas a \textbf{que se promuevan
nuevos talentos} / con vistas a \textbf{la promoción de nuevos talentos}}.

\begin{ejercicio}\label{ej:13-analitica}
A partir de las siguientes oraciones compuestas, describa analíticamente las
conjunciones subordinantes y, donde corresponda, los valores que manifiestan los
adverbios, preposiciones o locuciones que las preceden, según se muestra en el
ejemplo resuelto~1.

\begin{exe}
  \ex Actuó así a pesar de que le aconsejaron lo contrario.
  \ex Aunque llueva, vamos a salir.
  \ex Habló con el rector con el fin de que se aprobase su curso.
  \ex Lo explico de nuevo para que no queden dudas.
  \ex Me alegro porque lo hice con cariño.
  \ex No pude llegar antes porque había mucho tráfico.
  \ex No va a aceptar, aunque se lo pidamos amablemente.
  \ex Nos alegra tu visita porque te extrañábamos.
  \ex Te lo repito para que lo entiendas.
  \ex Trabajamos con el fin de que el proyecto sea un éxito.
\end{exe}

\begin{resuelto}
1. [\etq{Ocomp} Actuó así [\etq{SP} \textit{a pesar de} [\etq{OC}
\textbf{\textit{que}} le aconsejaron lo contrario]]], con valor concesivo en
\forma{a pesar de} (locución preposicional).
\end{resuelto}
\end{ejercicio}

\paragraph{Los relativos.}\label{sec:13-relativos}
Al igual que las conjunciones subordinantes, los relativos son los núcleos
funcionales de la oración (es decir, del sintagma complementante, la máxima
proyección oracional) y determinan la naturaleza de la oración subordinada como
relativa. El carácter pronominal de estas unidades es el factor clave que
establece una diferencia frente a las conjunciones subordinantes: dado que los
relativos constituyen expresiones referenciales, cumplen funciones sintácticas,
tanto argumentales como no argumentales, dentro de la oración que encabezan
(tabla~\ref{tab:13-relativos}).

{\small
\begin{longtable}{@{}>{\raggedright\arraybackslash}p{40mm}
    >{\raggedright\arraybackslash}p{83mm}@{}}
  \caption{Clases de relativos}\label{tab:13-relativos}\\
  \toprule
  \endfirsthead
  \toprule
  \endhead
  Con antecedente léxico explícito & Adverbios: \textit{como, cuando, donde,
    cuanto}. Determinantes: \textit{cuya / cuyo / cuyas / cuyos}. Pronombres
    (simples y complejos): \textit{que; quien / quienes; la / el cual, las / los
    cuales; el / la / los / las que}. \\
  \midrule
  Sin antecedente léxico explícito & Libres (antecedente implícito):
    \textit{como, cuando, cuanto, donde; quien / quienes; cuanta / cuanto /
    cuantas / cuantos}. Semilibres (antecedente tácito o elidido): \textit{el /
    la / los / las que, lo que}. \\
  \bottomrule
\end{longtable}}

Los relativos pueden tener un antecedente léxico explícito con el cual establecen
una relación de correferencia, como en \forma{El día que me quieras}, en el que
el pronombre relativo \forma{que} toma la referencia de su antecedente nominal
\forma{día}. Pero los relativos también pueden emplearse sin un antecedente
léxico explícito: en esta función se los denomina \termino{relativos libres}. En
su significado, estos relativos incorporan una palabra general como antecedente
implícito, como en \forma{Cuando tallan los recuerdos}, donde el relativo
\forma{cuando} se interpreta como \sig{en el momento en que}.

Los relativos simples (\forma{que, quien, cuanto, donde, cuando\ldots}) y los que
incorporan el artículo definido (\forma{el que, las que, lo que, los cuales, la
cual, lo cual\ldots}) presentan comportamientos morfológicos distintos: unos son
invariables (\forma{que, como, cuando, donde}), otros flexionan en número
(\forma{quien, quienes; cual, cuales}) y otros varían en género y número
(\forma{cuanto, cuanta, cuantos, cuantas; cuya, cuyo, cuyas, cuyos}).

\begin{recuadro}[Para ampliar]
La mayoría de los relativos son formas simples (una sola palabra), ya se trate de
\textit{relativos con antecedente léxico explícito}, como en
\exref{ex:13-ambulancia-a} y \exref{ex:13-ambulancia-b}, o de \textit{relativos
libres}, como en \exref{ex:13-ambulancia-c} y \exref{ex:13-ambulancia-d}:

\begin{exe}
  \ex Formas relativas simples (una sola palabra):
    \begin{xlist}
      \ex\label{ex:13-ambulancia-a} La ambulancia [\etq{OR} \textit{que}
          robaron los asaltantes] llevaba un paciente.
      \ex\label{ex:13-ambulancia-b} La ambulancia [\etq{OR} \textit{donde}
          estaba el paciente] fue robada por los asaltantes.
      \ex\label{ex:13-ambulancia-c} [\etq{OR} \textit{Quienes} robaron la
          ambulancia] habían realizado otros asaltos previamente.
      \ex\label{ex:13-ambulancia-d} Los asaltantes abandonaron la ambulancia
          [\etq{OR} por \textit{donde} la habían robado].
    \end{xlist}
\end{exe}

En \exref{ex:13-ambulancia-a} y \exref{ex:13-ambulancia-b}, \forma{que} y
\forma{donde} son relativos ligados anafóricamente a un antecedente léxico
(\forma{ambulancia}): \forma{que robaron\ldots} = \sig{los asaltantes robaron la
ambulancia}; \forma{donde estaba\ldots} = \sig{en la ambulancia estaba un
paciente}. En \exref{ex:13-ambulancia-c} y \exref{ex:13-ambulancia-d},
\forma{quienes} y \forma{donde} son relativos libres, sin un antecedente léxico
explícito: \forma{quienes robaron\ldots} = \sig{las personas que robaron la
ambulancia}; \forma{por donde la habían\ldots} = \sig{por el lugar donde la
habían robado} (antecedentes implícitos).

Adicionalmente, también hay formas relativas compuestas: los pronombres relativos
complejos, los pronombres relativos semilibres y los pronombres relativos con
antecedente elidido \parencite[§\,44.1d--g]{rae2025}. Los pronombres
\textit{relativos complejos} de \exref{ex:13-compuestas-a} y
\exref{ex:13-compuestas-b}, ligados anafóricamente a un antecedente léxico
explícito, forman unidades que no admiten material léxico intermedio. En
contraste, los pronombres \textit{relativos semilibres} de
\exref{ex:13-compuestas-c} y \exref{ex:13-compuestas-d}, con antecedente tácito
independiente del contexto, y los pronombres \textit{relativos con elisión del
antecedente} de \exref{ex:13-compuestas-e} y \exref{ex:13-compuestas-f}, con
antecedente dependiente del contexto, carecen de un antecedente explícito y
admiten material léxico intermedio, ya que se constituyen de dos unidades
diferentes:

\begin{exe}
  \ex Formas relativas compuestas (más de una palabra):
    \begin{xlist}
      \ex\label{ex:13-compuestas-a} La ambulancia [\etq{OR} en \textit{la que}
          huyeron los asaltantes] fue robada a los recolectores.
      \ex[*]{\label{ex:13-compuestas-b} La ambulancia [en \textit{la única que}
          huyeron los asaltantes] fue robada a los recolectores.}
      \ex\label{ex:13-compuestas-c} En realidad, la Ø [\etq{OR} \textit{que} se
          acordó de tu cumpleaños] fue mi hermana.
      \ex\label{ex:13-compuestas-d} En realidad, la única Ø [\etq{OR}
          \textit{que} se acordó de tu cumpleaños] fue mi hermana.
      \ex\label{ex:13-compuestas-e} Había dos ambulancias. Los asaltantes
          huyeron en la Ø [\etq{OR} \textit{que} llevaba un paciente].
      \ex\label{ex:13-compuestas-f} Había dos ambulancias. Los asaltantes
          huyeron en la única Ø [\etq{OR} \textit{que} llevaba un paciente].
    \end{xlist}
\end{exe}

En \exref{ex:13-compuestas-a}, \forma{la que} (= \forma{la cual}) es un pronombre
relativo complejo con antecedente léxico explícito: \forma{en la que
huyeron\ldots} = \sig{los asaltantes huyeron en la ambulancia}. En
\exref{ex:13-compuestas-c}, \forma{la que} es un pronombre relativo semilibre,
sin antecedente léxico explícito: \forma{la \{Ø\} que se acordó\ldots} = \sig{la
persona que se acordó} (antecedente tácito no contextual). En
\exref{ex:13-compuestas-e}, \forma{la que} es un pronombre relativo con elisión
de un antecedente léxico: \forma{la \{Ø\} que llevaba\ldots} = \sig{la
ambulancia que llevaba} (antecedente elidido contextual).
\end{recuadro}

\subsection{Tipo de flexión de las subordinadas}\label{sec:13-flexion}

Las oraciones subordinadas pueden articularse alrededor de formas verbales
finitas (verbos con flexión temporal-aspectual-modal y rasgos de concordancia
personal) o alrededor de formas verbales no finitas (verbos sin flexión
temporal-modal ni rasgos de concordancia personal), como el infinitivo, el
gerundio y el participio, diferenciados entre sí por su forma y sus valores
aspectuales (véase el \capref{11}).

Por un lado, las oraciones subordinadas con encabezadores léxicos (conjunciones
subordinantes o relativos) tienen asociadas flexiones finitas, cuyas
posibilidades témporo-modales dependen tanto de las correlaciones establecidas
con los verbos de la oración matriz como de los encabezadores que las introducen,
ya que las conjunciones subordinantes son más selectivas de los modos verbales
que los relativos.

\begin{exe}
  \ex\label{ex:13-modo1}
    \begin{xlist}
      \ex\label{ex:13-modo1-a} Xi Jinping se pregunta [\etq{OC} \textit{si} China
          tendrá nuevos éxitos económicos].
      \ex[*]{\label{ex:13-modo1-b} Xi Jinping se pregunta [\textit{si} China
          tenga nuevos éxitos económicos].}
    \end{xlist}
  \ex\label{ex:13-modo2}
    \begin{xlist}
      \ex\label{ex:13-modo2-a} Por eso, [\etq{SAdv} siempre [\etq{OC}
          \textit{que} podamos]], hagamos bien a todos.
      \ex[*]{\label{ex:13-modo2-b} Por eso, [siempre \textit{que} podemos],
          hagamos bien a todos.}
    \end{xlist}
  \ex\label{ex:13-modo3}
    \begin{xlist}
      \ex\label{ex:13-modo3-a} Busco a una persona [\etq{OR} \textit{que} sabe
          hablar castellano].
      \ex\label{ex:13-modo3-b} Busco una persona [\etq{OR} \textit{que} sepa
          hablar castellano].
    \end{xlist}
\end{exe}

El contraste de gramaticalidad entre \exref{ex:13-modo1-a} y
\exref{ex:13-modo1-b} revela que el modo indicativo depende de la conjunción
subordinante (\forma{si}), que marca la modalidad interrogativa de la oración
subordinada (dado que el modo subjuntivo es incompatible con la modalidad
interrogativa). La misma diferencia de aceptabilidad se advierte en el contraste
entre \exref{ex:13-modo2-a} y \exref{ex:13-modo2-b}, donde la lectura
condicional\footnote{La lectura temporal de \forma{siempre que}, como en
\forma{Siempre que podemos hacemos el bien a todos}, aparece asociada al modo
indicativo.} del sintagma encabezado por el adverbio \forma{siempre} se encuentra
asociada al modo subjuntivo de la subordinada encabezada por la conjunción
\forma{que}. En cambio, la buena formación de \exref{ex:13-modo3-a} y
\exref{ex:13-modo3-b} muestra que, en estos casos, el modo verbal indicativo o
subjuntivo no depende del pronombre relativo \forma{que}, sino del carácter
definido y existente (o no) de la entidad referida por el nombre.

Por otro lado, las oraciones subordinadas sin encabezadores léxicos, como las
interrogativas, se construyen con formas verbales finitas y no finitas, además de
ser compatibles con una conjunción subordinante
(\forma{que}).\footnote{Como ilustración histórica, ya en obras clásicas del Siglo
de Oro se evidencia que los ítems interrogativos constituyen focos
interrogativos y no núcleos de la máxima proyección oracional (SC), dado que
pueden coexistir con el núcleo \forma{que}: \forma{el soldado y yo le
preguntamos que quién era aquel caballero} (F. de Quevedo, \textit{Historia de
la vida del Buscón}). Iguales en su forma, los relativos, en cambio, no coexisten
con la conjunción subordinante. En la actualidad esta característica se puede
encontrar en la lengua coloquial.}

\begin{exe}
  \ex\label{ex:13-interr}
    \begin{xlist}
      \ex\label{ex:13-interr-a} Xi Jinping se pregunta [\etq{OC} qué nuevos
          éxitos económicos producirá China].
      \ex\label{ex:13-interr-b} Xi Jinping se pregunta [\etq{OC} cómo producir
          nuevos éxitos económicos en China].
    \end{xlist}
  \ex\label{ex:13-excl}
    \begin{xlist}
      \ex\label{ex:13-excl-a} No sabés [\etq{OC} cuánto me alegra ver que estás
          bien].
      \ex\label{ex:13-excl-b} Ya vas a ver [\etq{OC} qué impresionante es la
          casona que está en venta].
    \end{xlist}
\end{exe}

Los ejemplos \exref{ex:13-interr-a} y \exref{ex:13-interr-b} muestran la
compatibilidad de la flexión finita y no finita con los ítems interrogativos. En
cambio, las subordinadas que contienen ítems exclamativos, \exref{ex:13-excl-a} y
\exref{ex:13-excl-b}, considerados como focos, solo se construyen con una forma
verbal finita. El esquema de \exref{ex:13-diag7} muestra la interpretación de las
subordinadas con foco interrogativo y exclamativo.

\begin{exe}
  \ex\label{ex:13-diag7}
    \begin{xlist}
      \ex Foco interrogativo:\\[2pt]
          \begin{forest} gramatica
            [{SC = OC}
              [SN\\Foco [\textit{qué nuevos éxitos económicos}, roof]]
              [C\prima
                [C\\Ø]
                [SF [\textit{producirá China}, roof]]]]
          \end{forest}
      \ex Foco exclamativo:\\[2pt]
          \begin{forest} gramatica
            [{SC = OC}
              [SAdv\\Foco [\textit{cuánto}]]
              [C\prima
                [C\\Ø]
                [SF [\textit{me alegra ver que estás bien}, roof]]]]
          \end{forest}
    \end{xlist}
\end{exe}

\pendiente{Diagrama 7 sin imagen en el original; reconstruido}

En oraciones como \forma{Solicito se me inscriba fuera de término en condición de
alumno regular} o \forma{Espero me concedan una extensión del plazo
establecido}, normalmente asociadas a registros formales, se puede reponer la
conjunción subordinante \forma{que} sin alterar el significado oracional:
\forma{Solicito \textbf{que} se me inscriba fuera de término en condición de
alumno regular} o \forma{Espero \textbf{que} me concedan una extensión del plazo
establecido}. La tabla~\ref{tab:13-tipos} sintetiza los tipos de subordinadas
según su encabezador y su flexión.

{\small
\begin{longtable}{@{}>{\raggedright\arraybackslash}p{40mm}
    >{\raggedright\arraybackslash}p{83mm}@{}}
  \caption{Tipos de subordinadas según el encabezador y la flexión}\label{tab:13-tipos}\\
  \toprule
  \endfirsthead
  \toprule
  \endhead
  Subordinación con encabezadores léxicos explícitos &
    \textbf{Oraciones subordinadas completivas o conjuntivas (OC).} Conjunciones
    subordinantes de naturaleza conjuntiva encabezan el sintagma complementante
    (SC) sin cumplir una función sintáctica adicional. Formas verbales finitas
    (flexión temporal-aspectual-modal y rasgos de concordancia personal). \\
  \addlinespace
  & \textbf{Oraciones subordinadas relativas (OR).} Relativos de distinta
    naturaleza categorial (adverbios, determinantes o pronombres relativos)
    encabezan el sintagma complementante (SC) y, además, cumplen una función
    sintáctica nuclear. Formas verbales finitas. \\
  \midrule
  Subordinación sin encabezadores léxicos explícitos &
    \textbf{Oraciones subordinadas interrogativas y exclamativas indirectas
    (parciales).} Con foco interrogativo (adverbios, determinantes y pronombres
    interrogativos) y formas verbales finitas o no finitas. Con foco exclamativo
    (adverbios, determinantes y pronombres exclamativos) y formas verbales
    finitas. \\
  \addlinespace
  & \textbf{Oraciones subordinadas de \textit{que} facultativo.} Con potencial
    inserción de la conjunción y con formas verbales finitas (registros
    formales). \\
  \addlinespace
  & \textbf{Oraciones subordinadas de infinitivo, de gerundio y de participio.}
    Formas verbales no finitas. \\
  \bottomrule
\end{longtable}}

\begin{ejercicio}\label{ej:13-perifrasis}
En los siguientes ejemplos hay oraciones subordinadas con verbo flexionado y
otras con verbo no flexionado, además de predicados constituidos alrededor de una
perífrasis verbal (véase el \capref{11}). Según se muestra en los ejemplos
resueltos, distinga entre oraciones subordinadas (con o sin verbo en formas
personales) y perífrasis verbales y, donde corresponda, señale el elemento
introductor de la oración subordinada.

\begin{exe}
  \ex Como estoy cansado, no voy a la fiesta.
  \ex Estando cansado, no voy a la fiesta.
  \ex No voy a la fiesta porque estoy cansado.
  \ex Aun estando cansados, trabajamos duro.
  \ex Aunque estábamos cansados, trabajamos duro.
  \ex No estando invitado, no puedo ir al cumpleaños. Si no me invitan, no puedo
      ir al cumpleaños.
  \ex Es importante comprender el funcionamiento del motor.
  \ex Es importante que comprendas el funcionamiento del motor.
  \ex Logró que sus amigos supieran la verdad.
  \ex Logró saber la verdad.
  \ex Quiso que sus amigos supieran la verdad.
  \ex Quiso saber la verdad.
  \ex Se negó a que sus amigos supieran la verdad.
  \ex Se negó a saber la verdad.
  \ex De esta manera comenzó a saber la verdad.
  \ex De esta manera pudo saber la verdad.
  \ex De esta manera va a saber la verdad.
  \ex El tenor está cantando el aria.
  \ex El tenor ingresó a la escena mientras cantaba el aria.
  \ex El tenor ingresó cantando el aria.
\end{exe}

\begin{resuelto}[Ejemplos resueltos]
1. [\etq{Ocomp} [\etq{OS} \textbf{\textit{Como}} \textit{estoy cansado}], no voy a
la fiesta]: oración subordinada con verbo en forma personal introducida por la
conjunción subordinante \forma{como}. 2. [\etq{Ocomp} [\etq{OS} \textit{Estando
cansado}], no voy a la fiesta]: oración subordinada con verbo en forma no
personal, gerundio. 18. [\etq{O} El tenor \textit{está cantando} el aria]:
oración simple con perífrasis verbal de \forma{estar} + gerundio, de valor
progresivo, evento en curso.\pendiente{En 1 y 2 el original habla de
«perífrasis verbal», pero \textit{estoy cansado} y \textit{estando cansado} no
lo son}
\end{resuelto}
\end{ejercicio}

\section{Función sintáctica de las oraciones subordinadas}\label{sec:13-funcion}

\subsection{De la equivalencia funcional a una agrupación basada en la estructura interna}\label{sec:13-equivalencia}

Las oraciones subordinadas cumplen distintas funciones sintácticas en la oración
matriz de la que forman parte. Se trata de las mismas funciones sintácticas que
cumplirían un sintagma nominal, un sintagma adjetival, un sintagma adverbial o un
sintagma preposicional, aun cuando una oración subordinada es la máxima
proyección de una categoría funcional (sintagma complementante) y no de una
categoría léxica. Se comprende, entonces, que el criterio tradicional para
clasificar las oraciones subordinadas haya sido el de su \termino{equivalencia
funcional} con clases de palabras léxicas. Así, tanto en las gramáticas como en
los diseños curriculares, se suele distinguir, por su potencial funcional, entre
oraciones subordinadas sustantivas, adjetivas y adverbiales:

\begin{enumerate}[label=(\alph*)]
  \item Son \textit{oraciones subordinadas sustantivas} las que pueden cumplir
        las funciones sintácticas del sustantivo (sintagma nominal), por
        ejemplo, la de sujeto de la oración.
  \item Son \textit{oraciones subordinadas adjetivas} las que pueden cumplir las
        funciones sintácticas del adjetivo (sintagma adjetival), por ejemplo, la
        de modificador del nombre.
  \item Son \textit{oraciones subordinadas adverbiales} las que pueden cumplir
        las funciones sintácticas del adverbio (sintagma adverbial), por
        ejemplo, la de circunstancial en el predicado.
\end{enumerate}

Dejando de lado esta clasificación tradicional, en este capítulo, hasta ahora
hemos presentado las oraciones subordinadas en términos de su \termino{estructura
interna}. Nos hemos concentrado en la forma en que se articula la oración
subordinada (con encabezadores conjuntivos o relativos o sin encabezador) y en la
forma que adopta el verbo (flexionada o no finita). Teniendo en cuenta dichos
parámetros, podemos agrupar las oraciones subordinadas del siguiente modo:

\begin{enumerate}[label=(\alph*)]
  \item \termino{Oraciones subordinadas completivas} o \termino{conjuntivas}
        (OC), encabezadas por conjunciones subordinantes y con flexión verbal
        finita.
  \item \termino{Oraciones subordinadas relativas} (OR), encabezadas por
        pronombres relativos y con flexión verbal finita.
  \item \termino{Oraciones subordinadas sin encabezadores léxicos explícitos}
        (OS). Son las oraciones subordinadas interrogativas o exclamativas
        indirectas (con foco interrogativo o exclamativo), las oraciones
        subordinadas de \forma{que} facultativo (en las que se puede reponer la
        conjunción \forma{que}) y las oraciones subordinadas con formas verbales
        no finitas (infinitivo, gerundio, participio).
\end{enumerate}

\begin{recuadro}[Para ampliar]
Por fuera de las regularidades estructurales marcadas para las oraciones
subordinadas (con o sin encabezadores léxicos explícitos más flexión verbal), se
incluyen como subordinadas los siguientes tipos, que constituyen objetos
lingüísticos que no responden a estructuras oracionales como las que venimos
viendo hasta ahora:

\begin{enumerate}[label=(\alph*)]
  \item Las \textit{oraciones mencionadas}, que son las que normalmente se
        destacan en la escritura por medio de comillas (como en \forma{El
        ministro de Economía manifestó: «\textbf{El que apuesta al dólar
        pierde}»}) y pueden ser tanto oraciones completas como fragmentos
        oracionales (como en \forma{El ministro de Economía respondió: «En
        Argentina»}).
  \item Las \textit{fórmulas sintácticas}, que responden a esquemas fijos de
        distinto tipo (como en \forma{\textbf{Sea quien sea}, no contesto
        llamados}). Estas fórmulas pueden tener distintos formatos, alternativos
        y no predecibles (como en \forma{\textbf{Sea como sea}, ya terminaron las
        clases} y en \forma{\textbf{Como sea}, ya terminaron las clases}, o en
        \forma{\textbf{Con todo y estando enfermo}, asistió a clase} y en
        \forma{\textbf{Con todo}, asistió a clase}).
\end{enumerate}
\end{recuadro}

\subsection{Grado de integración de las subordinadas y funciones argumentales y no argumentales}\label{sec:13-integracion}

Las funciones sintácticas son dependencias estructurales que se establecen entre
un núcleo y un constituyente subordinado a él. En el marco de la oración simple,
se distingue entre funciones argumentales (exigidas por un núcleo léxico) y no
argumentales (compatibles con un núcleo, aunque no requeridas por él). Esta misma
caracterización se recoge en el marco de la oración compuesta, donde las
funciones sintácticas de las subordinadas están determinadas por el tipo de
relación con la oración matriz:

\begin{enumerate}[label=(\alph*)]
  \item La oración subordinada puede \textit{estar integrada} dentro de la
        oración matriz, donde cumple una función argumental o actúa como
        predicativo.
  \item La oración subordinada puede \textit{modificar} a la oración matriz,
        donde cumple una función no argumental.
\end{enumerate}

Importa diferenciar ambos tipos de relaciones (integración y modificación) para
comprender que, si bien la subordinación como vínculo interoracional es una, su
proyección en la estructura de la oración compuesta puede ser diversa. En unos
casos, la oración subordinada contribuye con la buena formación de la oración
matriz, que sin ella no estaría completa, mientras que en otros, su participación
dentro de la oración matriz es más periférica. Todas las oraciones subordinadas,
tanto las completivas como las relativas, pueden cumplir funciones sintácticas
argumentales y no argumentales.

En suma, las oraciones subordinadas cumplen funciones sintácticas que se pueden
agrupar por su aporte a la estructura oracional y por el nivel sintagmático en
el que se proyectan:

\begin{enumerate}[label=(\alph*)]
  \item \textit{Funciones sintácticas argumentales y predicativas}. Las oraciones
        subordinadas completivas y relativas sin antecedente léxico pueden
        cumplir funciones argumentales (sujeto de la oración, complemento del
        verbo, complemento o término de preposición) y también pueden ser
        predicativos (subjetivos u objetivos).
  \item \textit{Funciones sintácticas no argumentales}:
        \begin{enumerate}[label=(b.\arabic*)]
          \item \textit{Modificadoras del núcleo nominal o adverbial}. Las
                oraciones subordinadas relativas con antecedente léxico siempre
                cumplen funciones no argumentales de modificación del nombre o
                del adverbio (adjunto nominal o adverbial).
          \item \textit{Modificadoras del núcleo verbal}. Las oraciones
                subordinadas relativas sin antecedente léxico pueden cumplir
                funciones no argumentales respecto del verbo en el predicado
                (adjunto circunstancial).
          \item \textit{Modificadoras del núcleo oracional}. Las oraciones
                subordinadas completivas, complemento de preposición o de
                adverbio, pueden formar construcciones que modifican el
                contenido proposicional de la oración (adjunto modificador de
                núcleo oracional, MNO).
          \item \textit{Modificadoras de la modalidad oracional}. Las oraciones
                subordinadas completivas, complemento de preposición o de
                adverbio, pueden formar construcciones que modifican la
                modalidad (adjunto modificador de modalidad, MMO).
          \item \textit{Modificadoras del grado}. Las oraciones subordinadas
                completivas pueden cumplir funciones de modificación
                (complementación o adjunción, según la perspectiva que se adopte)
                correlativas de cuantificadores de intensificación en estructuras
                cuantificadas complejas (coda comparativa y coda consecutiva).
        \end{enumerate}
\end{enumerate}

\subsection{Funciones de las oraciones subordinadas completivas}\label{sec:13-completivas}

Las oraciones subordinadas completivas, que son introducidas por una conjunción
subordinante y contienen un verbo en forma finita, las oraciones subordinadas
interrogativas y exclamativas indirectas y las oraciones subordinadas de
infinitivo pueden cumplir funciones argumentales respecto del verbo (sujeto y
complemento objeto directo) o funciones argumentales respecto de la preposición
(complemento o término). Se trata de funciones características de los nombres o
sintagmas nominales, por lo que tradicionalmente se las denominó
\textit{sustantivas}.

\begin{exe}
  \ex\label{ex:13-suj} Sujeto de la oración (argumento verbal):
    \begin{xlist}
      \ex\label{ex:13-suj-a} Me disgusta [\etq{OC/S} que plantees problemas].
          (Me disgusta \textit{tu planteo}).
      \ex\label{ex:13-suj-b} Finalmente se conoce [\etq{OC/S} qué cifra alcanzó
          la inflación el mes pasado]. (Finalmente se conoce \textit{la cifra de
          inflación del mes pasado}).
      \ex\label{ex:13-suj-c} Es recomendable [\etq{OC/S} salir temprano para la
          costa]. (Es recomendable \textit{una salida temprana}).
    \end{xlist}
  \ex\label{ex:13-cod} Complemento objeto directo (argumento verbal):
    \begin{xlist}
      \ex\label{ex:13-cod-a} Estamos esperando [\etq{OC/COD} que plantees
          problemas]. (Estamos esperando \textit{tus planteos}).
      \ex\label{ex:13-cod-b} El Presidente desconoce [\etq{OC/COD} qué cifra
          alcanzó la inflación el mes pasado]. (El Presidente desconoce \textit{la
          cifra de inflación del mes pasado}).
      \ex\label{ex:13-cod-c} Ayer decidimos [\etq{OC/COD} salir temprano para la
          costa]. (Ayer decidimos \textit{una salida temprana}).
    \end{xlist}
  \ex\label{ex:13-ct} Complemento o término de preposición (argumento
      preposicional):
    \begin{xlist}
      \ex\label{ex:13-ct-a} Mi respuesta va a depender [de [\etq{OC/CT} si
          planteás problemas]]. (Mi respuesta va a depender de \textit{tus
          planteos}).
      \ex\label{ex:13-ct-b} Ahora dan explicaciones [sobre [\etq{OC/CT} cuál fue
          la inflación del mes pasado]]. (Ahora dan explicaciones sobre \textit{la
          inflación del mes pasado}).
      \ex\label{ex:13-ct-c} Al final optamos [por [\etq{OC/CT} salir temprano para
          la costa]]. (Al final optamos por \textit{una salida temprana}).
    \end{xlist}
\end{exe}

En las oraciones de (a), la subordinada introducida por una conjunción cumple las
funciones argumentales de sujeto \exref{ex:13-suj-a}, complemento objeto directo
\exref{ex:13-cod-a} y complemento o término de preposición \exref{ex:13-ct-a}.
Asimismo, en las oraciones de (b), la subordinada interrogativa indirecta cumple
las funciones argumentales de sujeto \exref{ex:13-suj-b}, complemento objeto
directo \exref{ex:13-cod-b} y complemento o término de preposición
\exref{ex:13-ct-b}. Finalmente, en las oraciones de (c), la subordinada de
infinitivo cumple las mismas funciones argumentales de sujeto
\exref{ex:13-suj-c}, complemento objeto directo \exref{ex:13-cod-c} y complemento
o término de preposición \exref{ex:13-ct-c}. En todos los casos se ilustra entre
paréntesis la equivalencia nominal de tales oraciones subordinadas.

Aunque las oraciones subordinadas, según se ve, pueden cumplir funciones
sintácticas idénticas a las de los nombres o sintagmas nominales, esto no
significa que se trate de sustantivos, ya que existen diferencias interpretativas
entre las oraciones subordinadas y los nombres. En este sentido, las
subordinadas completivas, por su valor modal, pueden ser enunciativas, como en
\exref{ex:13-enun-a}, \exref{ex:13-enun-b} y \exref{ex:13-enun-c}, exhortativas
\exref{ex:13-enun-d}, interrogativas \exref{ex:13-indir}, tanto
totales\footnote{Las interrogativas totales o disyuntivas se caracterizan por
presentar opciones entre las que se debe elegir: \forma{¿Querés ir al cine (o
no)?; ¿Querés ir al cine o al teatro?} Las interrogativas parciales o
pronominales van encabezadas por pronombres o adverbios interrogativos que son el
foco de la pregunta sobre el que se debe responder: \forma{¿Dónde querés ir esta
noche?; ¿Quién te lo dijo?}} \exref{ex:13-indir-a} como parciales,
\exref{ex:13-indir-b} y \exref{ex:13-indir-c}, o
exclamativas\footnote{\label{fn:13-exclamativas}Las oraciones exclamativas son
siempre parciales, ya que se inician con un ítem exclamativo: \forma{¡Qué
barbaridad!; ¡Cuánta injusticia!}} \exref{ex:13-indign}, según los verbos que las
seleccionan y los valores semánticos que aportan. Esto marca una diferencia con
los sintagmas nominales, ya que las subordinadas son seleccionadas en función de
sus valores modales; en cambio, los nombres, que carecen de estos valores, son
seleccionados en función de su semántica.

\begin{exe}
  \ex Oraciones subordinadas completivas enunciativas y exhortativas:
    \begin{xlist}
      \ex\label{ex:13-enun-a} El estudiante responde [\etq{OC/COD} que los
          ángulos internos de un triángulo suman 180°].
      \ex\label{ex:13-enun-b} El estudiante espera [\etq{OC/COD} que los ángulos
          internos de un triángulo sumen 180°].
      \ex\label{ex:13-enun-c} El estudiante desea [\etq{OC/COD} aprobar el
          examen de matemática].
      \ex\label{ex:13-enun-d} El estudiante solicita [\etq{OC/COD} que se
          prorroguen las clases].
    \end{xlist}
  \ex\label{ex:13-indir} Oraciones subordinadas completivas interrogativas
      indirectas (totales y parciales):
    \begin{xlist}
      \ex\label{ex:13-indir-a} El maestro pregunta [\etq{OC/COD} si los ángulos
          internos de un triángulo suman 180°].
      \ex\label{ex:13-indir-b} El maestro pregunta [\etq{OC/COD} cuánto suman los
          ángulos internos de un triángulo].
      \ex\label{ex:13-indir-c} El maestro desconoce [\etq{OC/COD} hasta cuándo se
          prorrogarán las clases].
    \end{xlist}
  \ex\label{ex:13-indign} Oración subordinada completiva exclamativa indirecta
      (parcial): Es indignante [\etq{OC/S} cuánto demora el trámite de título].
\end{exe}

Las oraciones subordinadas completivas también pueden cumplir la función de
predicativo subjetivo \exref{ex:13-ps} u objetivo \exref{ex:13-po}.

\begin{exe}
  \ex\label{ex:13-ps} Predicativo subjetivo (predicado semántico):
    \begin{xlist}
      \ex Mi esperanza es [\etq{OC/PS} que no surjan problemas].
      \ex La duda es [\etq{OC/PS} qué pasa con el feriado del mes próximo].
      \ex El objetivo de la política fiscal es [\etq{OC/PS} eliminar los
          impuestos distorsivos].
    \end{xlist}
  \ex\label{ex:13-po} Predicativo objetivo (predicado semántico):
    \begin{xlist}
      \ex Ayer vi a tu hermano [\etq{OC/PO} entrar a la oficina]. (Ayer vi
          \textit{que tu hermano entraba a la oficina}).
      \ex En la esquina hay una cuadrilla de operarios [\etq{OC/PO} rompiendo la
          vereda].
    \end{xlist}
\end{exe}

\subsection{Funciones de las oraciones subordinadas relativas}\label{sec:13-funcrelativas}

Las oraciones subordinadas relativas (encabezadas por un relativo) pueden
construirse como ligadas a un antecedente léxico explícito de naturaleza nominal
o bien sin antecedente, caso en que se denominan \textit{relativas libres}.

\subsubsection{Relativas con antecedente}\label{sec:13-conantecedente}

Las relativas ligadas a un antecedente pueden cumplir funciones no argumentales
como modificadoras del nombre en calidad de adjuntos, que pueden ser
\termino{especificativos} o \termino{explicativos}. Los adjuntos especificativos,
como los de \exref{ex:13-as}, restringen la semántica del nombre al que
modifican; en cambio, los adjuntos explicativos, como los de \exref{ex:13-ax},
amplían y detallan la semántica del sintagma nominal al que se adjuntan. En ambos
casos, las funciones que cumplen son características de los adjetivos o
sintagmas adjetivales.

\begin{exe}
  \ex\label{ex:13-as} Adjunto nominal especificativo:
    \begin{xlist}
      \ex\label{ex:13-as-a} La policía detuvo a dos delincuentes [\etq{OR/AS} que
          atracaron un banco]. (La policía detuvo a dos delincuentes. \textit{Dos
          delincuentes} atracaron un banco).
      \ex\label{ex:13-as-b} Las personas [\etq{OR/AS} a quienes se detuvo]
          negaron su responsabilidad en el hecho. (Las personas negaron su
          responsabilidad en el hecho. Se detuvo \textit{a esas personas}).
      \ex\label{ex:13-as-c} Dijeron no conocer el banco [\etq{OR/AS} donde se
          produjo el atraco]. (Dijeron no conocer el banco. \textit{En el banco}
          se produjo el atraco).
      \ex\label{ex:13-as-d} Añadieron que para la fecha [\etq{OR/AS} cuando se
          produjo el hecho] estaban fuera de la ciudad. (Añadieron que para la
          fecha estaban fuera de la ciudad. \textit{En la fecha} se produjo el
          hecho).
    \end{xlist}
  \ex\label{ex:13-ax} Adjunto nominal explicativo:
    \begin{xlist}
      \ex\label{ex:13-ax-a} La causa, [\etq{OR/AX} la cual quedó a cargo del
          juez], suma testimonios de varios testigos. (La causa suma testimonios
          de varios testigos. \textit{La causa} quedó a cargo del juez).
      \ex\label{ex:13-ax-b} Según un testigo, [\etq{OR/AX} cuyo testimonio se
          considera clave], los ladrones fueron vistos en el banco. (Según un
          testigo, los ladrones fueron vistos en el banco. \textit{Su}
          testimonio se considera clave).
    \end{xlist}
\end{exe}

En ambas series, el relativo que encabeza la oración subordinada puede tener
funciones nominales, como en \exref{ex:13-as-a}, \exref{ex:13-as-b} y
\exref{ex:13-ax-a}, adverbiales, como en \exref{ex:13-as-c} y
\exref{ex:13-as-d}, o de determinante, como en \exref{ex:13-ax-b}. En todos los
casos, el relativo liga su interpretación al antecedente nominal (es decir, se
carga de significado a partir del nombre o sintagma nominal que lo antecede).

\subsubsection{Relativas libres}\label{sec:13-libres}

Las oraciones subordinadas relativas libres, encabezadas por un relativo no
ligado a un antecedente léxico explícito, pueden cumplir funciones nominales o
adverbiales.

\paragraph{Relativas libres con funciones adverbiales.}\label{sec:13-libres-adv}
Las relativas libres encabezadas por los adverbios \forma{donde, cuanto, cuando}
y \forma{como} pueden desempeñarse, dentro del predicado oracional, como
circunstanciales, argumentales o no argumentales, respecto del verbo, como se
presenta en \exref{ex:13-circ}. Las funciones que cumplen son características de
los sintagmas adverbiales y preposicionales.

\begin{exe}
  \ex\label{ex:13-circ} Circunstancial (complemento o adjunto verbal):
    \begin{xlist}
      \ex\label{ex:13-circ-a} Corrió en la cinta [\etq{OR/AC} cuanto le dieron
          sus piernas]. (Cantidad).
      \ex\label{ex:13-circ-b} Terminó la primaria [\etq{OR/AC} donde la había
          empezado]. (Lugar).
      \ex\label{ex:13-circ-c} Habla inglés [\etq{OR/AC} como le enseñaron en la
          casa]. (Modo o manera).
      \ex\label{ex:13-circ-d} Salió de su casa [\etq{OR/AC} cuando lo llamaron del
          campo]. (Tiempo).
    \end{xlist}
\end{exe}

Las oraciones de \exref{ex:13-circ} se corresponden con las relativas libres que,
en la tradición gramatical, fueron denominadas \textit{subordinadas adverbiales
propias} porque podían ser reemplazadas por adverbios. Así, podemos ver que las
relativas libres funcionan como adverbios: en \exref{ex:13-circ-a}, como adjunto
circunstancial de cantidad, encabezada por el relativo \forma{cuanto}; en
\exref{ex:13-circ-b}, como adjunto circunstancial de lugar, encabezada por el
relativo \forma{donde}; en \exref{ex:13-circ-c}, como adjunto circunstancial de
modo, encabezada por el relativo \forma{como}; y en \exref{ex:13-circ-d}, como
adjunto circunstancial de tiempo, encabezada por el relativo
\forma{cuando}.\footnote{En este apartado solo consideramos las subordinadas
relativas introducidas por \forma{cuando}; el resto de las construcciones que
expresan \sig{tiempo} se ven en §\ref{sec:13-temporales}.}

Es de notar que los mismos adverbios relativos libres de los ejemplos de
\exref{ex:13-circ} pueden estar ligados a un antecedente léxico explícito
(\forma{todo, escuela, manera, momento}), como en \exref{ex:13-circant}. En estos
casos la construcción subordinada se convierte en una relativa con antecedente
que funciona como adjunto especificativo. Asimismo, se advierte una
correspondencia con una serie de adverbios no relativos (\forma{mucho, allí, así,
entonces}), funcionalmente equivalentes a las respectivas oraciones subordinadas
relativas libres.

\begin{exe}
  \ex\label{ex:13-circant} Relativas con antecedente en estructuras con función
      de circunstancial:
    \begin{xlist}
      \ex Corrió en la cinta todo \textit{cuanto le dieron sus piernas}. Corrió
          \textit{mucho} en la cinta.
      \ex Terminó la primaria en la escuela \textit{donde la había empezado}. La
          terminó \textit{allí}.
      \ex Habla inglés de la manera \textit{como le enseñaron en la casa}. Lo
          habla \textit{así}.
      \ex Salió de su casa en el momento \textit{cuando lo llamaron del campo}.
          Salió \textit{entonces}.
    \end{xlist}
\end{exe}

\paragraph{Relativas libres con funciones nominales.}\label{sec:13-libres-nom}
Las relativas libres (encabezadas por el pronombre relativo \forma{quien} y su
variante plural) y semilibres (encabezadas por el pronombre relativo \forma{el
que} y sus variantes de género y número) pueden cumplir funciones argumentales
características de los nombres o sintagmas nominales, relacionadas con el verbo
(sujeto y complemento objeto directo) y la preposición (complemento o término).

\begin{exe}
  \ex\label{ex:13-abarca-s} Sujeto de la oración (argumento verbal):
    \begin{xlist}
      \ex\label{ex:13-abarca-s-a} {}[\etq{OR/S} Quien mucho abarca], poco
          aprieta. (\textit{La persona que mucho abarca}, poco aprieta).
      \ex\label{ex:13-abarca-s-b} {}[\etq{OR/S} Los que mucho abarcan], poco
          aprietan. (\textit{Los hombres que mucho abarcan}, poco aprietan).
    \end{xlist}
  \ex\label{ex:13-abarca-cod} Complemento objeto directo (argumento verbal):
    \begin{xlist}
      \ex\label{ex:13-abarca-cod-a} A [\etq{OR/COD} quien mucho abarca] en nada
          se lo considera idóneo. (A \textit{la persona que mucho abarca} en nada
          se la considera idónea).
      \ex\label{ex:13-abarca-cod-b} A [\etq{OR/COD} las que mucho abarcan] en nada
          se las considera idóneas. (A \textit{las personas que mucho abarcan} en
          nada se las considera idóneas).
    \end{xlist}
  \ex\label{ex:13-abarca-ct} Complemento o término de preposición (argumento
      preposicional):
    \begin{xlist}
      \ex\label{ex:13-abarca-ct-a} Todos se alejan [de [\etq{OR/CT} quienes mucho
          abarcan]]. (Todos se alejan de \textit{las personas que mucho
          abarcan}).
      \ex\label{ex:13-abarca-ct-b} Todos se alejan [de [\etq{OR/CT} la que mucho
          abarca]]. (Todos se alejan de \textit{la persona que mucho abarca}).
    \end{xlist}
\end{exe}

En las oraciones (a), la relativa libre cumple las funciones argumentales de
sujeto \exref{ex:13-abarca-s-a}, complemento objeto directo
\exref{ex:13-abarca-cod-a} y complemento o término de preposición
\exref{ex:13-abarca-ct-a}. El mismo esquema de funciones sintácticas se replica
con la relativa semilibre en la serie (b): sujeto \exref{ex:13-abarca-s-b},
complemento objeto directo \exref{ex:13-abarca-cod-b} y complemento o término de
preposición \exref{ex:13-abarca-ct-b}. En todos estos casos, las relativas
muestran una equivalencia funcional con sintagmas nominales encabezados por un
nombre general con rasgo [+humano] (\forma{hombre, mujer, persona}) o con
pronombres personales (\forma{él, ella, ellas, ellos}), hecho que normalmente
deriva en su consideración como \termino{relativas sustantivadas}.

En relación con su valor sustantivado conviene señalar que, aun siendo todas
oraciones subordinadas, existen diferencias interpretativas y sintácticas,
además de diferencias estructurales, entre las completivas y las relativas
sustantivadas. Así, las oraciones subordinadas completivas (OC) denotan nociones
abstractas (contenidos proposicionales) y se asimilan a oraciones
\exref{ex:13-esto}; en cambio, las relativas (OR) libres, aunque también pueden
cumplir funciones argumentales, aluden a cualquier referente y se asimilan a
nombres o pronombres \exref{ex:13-este}:

\begin{exe}
  \ex\label{ex:13-esto} Oraciones subordinadas completivas (OC), asimilables a
      oraciones, conmutables por el demostrativo neutro:
    \begin{xlist}
      \ex Me disgusta [\etq{OC/S} que plantees problemas]. Me disgusta
          \textit{esto}. (Pero *Me disgusta \textit{este / él}).
      \ex Finalmente se conoce [\etq{OC/S} qué cifra alcanzó la inflación el mes
          pasado]. Finalmente se conoce \textit{esto}. (Pero *Finalmente se conoce
          \textit{esta / ella}).
      \ex Es recomendable [\etq{OC/S} salir temprano para la costa]. Es
          recomendable \textit{esto}. (Pero *Es recomendable \textit{este / él}).
    \end{xlist}
  \ex\label{ex:13-este} Oraciones subordinadas relativas (OR), asimilables a
      nombres o pronombres, conmutables por pronombres demostrativos o
      personales:
    \begin{xlist}
      \ex {}[\etq{OR/S} Quien mal anda], mal acaba. \textit{Este / Él} mal acaba.
          (Pero *\textit{Esto} mal acaba). \textit{La persona que mal anda} mal
          acaba.
      \ex {}[\etq{OR/S} Lo que mal empieza], mal acaba. \textit{Esto} mal acaba.
          (Pero *\textit{Este / Él} mal acaba). \textit{Aquello que mal empieza}
          mal acaba.
    \end{xlist}
\end{exe}

Las mismas oraciones subordinadas relativas libres y semilibres (pronombres
relativos \forma{quien} y \forma{el que}, con sus variantes) pueden cumplir la
función de predicativo subjetivo.

\begin{exe}
  \ex\label{ex:13-predrel} Predicativo subjetivo (predicado semántico):
    \begin{xlist}
      \ex\label{ex:13-predrel-a} ¿Usted es [\etq{OR/PS} quien siempre avisa de
          los cortes de servicio a EDESUR]? (¿Usted es \textit{la persona que
          siempre avisa de los cortes de servicio a EDESUR}?)
      \ex\label{ex:13-predrel-b} Yo soy [\etq{OR/PS} la que lleva esta casa
          adelante]. (Yo soy \textit{la mujer que lleva esta casa adelante}).
    \end{xlist}
\end{exe}

En las oraciones de \exref{ex:13-predrel}, la relativa libre
\exref{ex:13-predrel-a} y la semilibre \exref{ex:13-predrel-b} cumplen la
función sintáctica de predicativo subjetivo del verbo copulativo \forma{ser}. El
pronombre relativo replica la referencia presente en el contexto oracional
cuando el verbo de la relativa muestra una concordancia diferente de la tercera
persona, como en \forma{Yo soy el que soy}.

\begin{ejercicio}\label{ej:13-funciones}
Volviendo sobre las paremias del ejercicio~\ref{ej:13-paremias}, identifique las
funciones sintácticas que cumplen las oraciones subordinadas reconocidas en
ellas, según se muestra en los ejemplos resueltos.

\begin{resuelto}[Ejemplos resueltos]
1. [\etq{Ocomp} A la cama no te irás \textit{sin} [\etq{OC/CT} \textit{saber una
cosa más}]]: la OC de infinitivo (sin subordinante) es el complemento o término
(CT) de la preposición \forma{sin} (valor de polarización negativa: \sig{no
sabiendo una cosa más}). 2. [\etq{Ocomp} A [\etq{OR/COD} \textit{quien
madruga}], Dios lo ayuda]: la OR libre encabezada por el relativo \forma{quien}
(\sig{alguien madruga}) es el complemento objeto directo (COD) del verbo
transitivo \forma{ayudar}.
\end{resuelto}
\end{ejercicio}

\begin{ejercicio}\label{ej:13-relativos}
En las siguientes oraciones compuestas, según se muestra en los ejemplos
resueltos, identifique el relativo y la función sintáctica que cumple dentro de
la oración subordinada y, donde corresponda, señale el antecedente del relativo.

\begin{exe}
  \ex Ayudemos a quienes lo necesiten.
  \ex Compré una valija. De hecho, compré la que me recomendaste.
  \ex Conocí a la persona para quien trabajó mi padre.
  \ex Decime lo que pensás.
  \ex Dieron un premio a quienes participaron.
  \ex El coche que está estacionado allí es el mío.
  \ex El libro que me prestaste es muy interesante.
  \ex El que llegó tarde no entró.
  \ex Esa es la casa cuya puerta estaba pintada de verde.
  \ex Esta es la lapicera con la cual escribí la carta.
  \ex Hizo la preparación como le enseñaron.
  \ex La docente con quien estudié inglés era británica.
  \ex La película por la cual vinimos al festival fue excelente.
  \ex Lo voy a comprar cuando esté en oferta.
  \ex Los estudiantes que estudiaron aprobaron el examen.
  \ex Los estudiantes, que estudiaron, aprobaron el examen.
  \ex Necesito lo que preparaste para el relleno de la tarta.
  \ex No me importa lo que digan.
  \ex Quien persevera, vence.
  \ex Quienes estudian, aprenden.
\end{exe}

\begin{resuelto}[Ejemplos resueltos]
1. [\etq{Ocomp} Ayudemos a [\etq{OR/COD} \textit{quienes} lo necesiten]]:
\forma{quienes}, sujeto de la OR; sin antecedente. 3. [\etq{Ocomp} Conocí a la
persona [\etq{OR/AS} para \textit{quien} trabajó mi padre]]: \forma{quien},
complemento o término de la preposición \forma{para}; antecedente,
\forma{persona}.\pendiente{El original atribuye a \textit{para} «valor
sociativo»}
\end{resuelto}
\end{ejercicio}

\section{Las construcciones con valores adverbiales}\label{sec:13-adverbiales}

Como ya hemos adelantado, partiendo del supuesto de la equivalencia funcional con
el adverbio, la gramática tradicional agrupaba, bajo la etiqueta de
\textit{subordinadas adverbiales}, construcciones sintáctica y semánticamente muy
diversas. En este capítulo, en cambio, hemos privilegiado un enfoque que parte de
la conformación interna de las subordinadas. Por tanto, las consideradas
tradicionalmente \textit{adverbiales propias} (temporales de \forma{cuando},
locativas, modales y de cantidad), que podían ser sustituidas por un adverbio de
significado equivalente, han sido incluidas entre las relativas libres por estar
encabezadas por adverbios relativos (respectivamente \forma{cuando, donde, como}
y \forma{cuanto}; véase §\ref{sec:13-libres}).

En este apartado nos vamos a ocupar de las temporales no relativas y de las que
recibían la denominación de \textit{adverbiales impropias} (causales, finales,
condicionales, concesivas, comparativas y consecutivas), a partir de que se
reconocía que no existen adverbios por los que puedan ser reemplazadas. En todos
estos casos, según el enfoque analítico que propiciamos, se trata de estructuras
encabezadas por conjunciones subordinantes (\forma{que, si}), que introducen
oraciones subordinadas completivas, que son seleccionadas por adverbios (p.~ej.,
\forma{aun, antes} o \forma{siempre}), preposiciones (p.~ej., \forma{para, por,
de}) y locuciones (\forma{dado que, puesto que}) o bien resultan correlativas de
elementos intensificativos presentes en la oración principal (\forma{tanto, más,
menos}). En este sentido, vamos a referirnos a todas ellas como
\termino{construcciones con valores adverbiales}, denominación amplia que
utilizamos para el conjunto formado por el adverbio, la preposición o locución
más la subordinada, que en conjunto forman un sintagma adverbial o
preposicional, según los casos. En cuanto a las comparativas y consecutivas, que
«están entre las que presentan mayor complejidad sintáctica de toda la gramática
española» \parencite[§\,45.1a]{rae2025}, la construcción abarca la estructura
formada por el intensificador, el núcleo al que este especifica y la subordinada
completiva introducida por \forma{que}: \forma{Esa foto tuya es [más actual que
la que está en el álbum]}. El valor semántico (temporal, causal, condicional,
comparativo, etc.) que toma la construcción lo aporta el adverbio, la
preposición, una locución o bien la correlación con el elemento intensificador.
Así, a modo de ejemplo, podemos mencionar casos como \forma{antes (de), después
(de), siempre} (temporal); \forma{por, a causa de, debido a} (causal); \forma{aun,
a pesar de, pese a} (concesivo); \forma{en caso de} (condicional); \forma{para, a
fin de, con vistas a} (final); \forma{más\ldots\ que\ldots} (comparativo);
\forma{tanto\ldots\ que} (consecutivo).

Antes de ver más en detalle los distintos tipos de construcciones con valores
adverbiales, vamos a presentar un enfoque que, desde la perspectiva semántica,
nos permite introducir una propuesta integradora que supere la atomización de
valores con que tradicionalmente se presentan estas subordinadas, a partir del
reconocimiento de las operaciones básicas que subyacen a las conexiones que
establecen estas construcciones.

\subsection{Una perspectiva integradora}\label{sec:13-integradora}

Partimos de la idea de que la oración constituye la representación de un evento
o estado de cosas del mundo exterior (como en \forma{Llueve}) o interior (como en
\forma{Tengo frío}). Cuando combinamos oraciones, esas representaciones se
entrelazan en la mente del hablante, lo que lleva a que las exprese combinadas
secuencialmente (como en \forma{Llueve; tengo frío}), ya sea en estructuras
independientes, coordinadas (como en \forma{Llueve y tengo frío}) o integradas
mediante subordinación (como en \forma{Tengo frío porque llueve}).

Según plantea \textcite[4]{duque2025}, «todos los hablantes comprendemos el
mundo y los textos mediante comparaciones, generalizaciones u ordenando la
información en función de relaciones temporales o causales». Estas relaciones
resultan intuitivamente simples y, por lo tanto, sumamente accesibles porque se
fundamentan en operaciones básicas. En esta perspectiva,
\textcite{rudolph1988,rudolph1996} propuso un esquema explicativo que permite dar
cuenta de las relaciones semánticas fundamentales que manifestamos al conectar
nuestras ideas mediante la combinación de oraciones. Para esta autora
\parencite[20]{rudolph1996}, existen cuatro formas básicas de conexión:

\begin{description}
  \item[Adición {[A + B]}:] dos proposiciones A y B están estrechamente
        relacionadas y son válidas simultáneamente.
  \item[Contraste {[A $-$ B]}:] para el hablante, dos proposiciones A y B son
        válidas simultáneamente y la proposición B marca un contraste con la
        proposición A.
  \item[Tiempo {[A, B]}:] dos proposiciones A y B no son válidas simultáneamente,
        y la proposición B denota un evento que sigue en el tiempo a la
        proposición A.
  \item[Causalidad {[A > B]}:] para el hablante, dos proposiciones A y B pueden
        ser válidas simultáneamente o no, y están ligadas en una cadena causal o
        de necesidad lógica o empírica.
\end{description}

Estas cuatro formas básicas de conexión se emplean para vincular eventos o
estados de cosas. Por un lado, las relaciones de adición y tiempo permiten al
hablante organizar la presentación de los hechos. Por otro lado, las relaciones
de contraste y causalidad hacen posible que, además de establecer un orden entre
los sucesos, el hablante también exprese sus opiniones y puntos de vista sobre
los hechos, estableciendo entre ellos semejanzas y diferencias o relaciones de
causa-efecto. Así, bajo la perspectiva aquí descripta, las construcciones que
manifiestan relaciones lógicas, como las temporales, condicionales, concesivas,
causales y finales, y las que expresan valores cuantificacionales-intensificativos,
como las comparativas y consecutivas, entre las que las descripciones
gramaticales tradicionales no solían establecer conexiones, se resuelven en
relaciones más básicas que van de la simple conexión y ordenación de los eventos
(simultaneidad o secuencia) a la manifestación de contrastes (oposiciones y
comparaciones) y diferentes tipos de causalidades (causa real, final, virtual e
ineficiente) entre sucesos.

Según plantea \textcite{rudolph1988}, los valores semánticos que expresan estas
cuatro formas básicas de conexión interoracional se pueden proyectar mediante la
simple secuenciación de oraciones independientes, es decir, sin utilizar
conectores léxicos. Sin embargo, el empleo de conjunciones y locuciones
conjuntivas que hagan explícitas las relaciones lógicas, discursivas y
argumentativas que contraen las oraciones constituye una excelente orientación
para el oyente en cuanto a la interpretación de los textos. En este sentido, los
conectores y conjunciones de argumentación (aditivos, adversativos, concesivos,
consecutivos), de reformulación (reformuladores, explicativos, rectificativos,
recapitulativos) y de estructuración (organizativos, ejemplificativos,
digresivos) presentados en el \capref{12}, §\ref{sec:12-marcadores}, funcionan
como formas lingüísticas especializadas en expresar esos señalamientos dirigidos
al receptor.

Las relaciones subordinativas (y también las coordinativas, vistas en el
\capref{12}), que involucran la combinación de unidades de nivel oracional, son
responsables de dar unidad y sentido a los textos estableciendo conexiones entre
las ideas. Desde esta perspectiva de integración, a continuación vamos a
presentar dos subtipos de construcciones con valores adverbiales: a) las que
manifiestan valores lógico-relacionales (temporales, causales, finales,
condicionales y concesivas) y b) las que manifiestan valores
cuantificacionales-intensificativos (comparativas y consecutivas).

\subsection{Construcciones con valores lógico-relacionales}\label{sec:13-logicas}

Según vimos, las oraciones subordinadas con funciones argumentales forman parte
de una estructura compleja, en la que completan el valor requerido por un núcleo
(verbal o preposicional) del que dependen. En cuanto a las oraciones relativas
ligadas a un antecedente nominal, contribuyen semánticamente a la restricción
(relativas especificativas) o a la explicación (relativas explicativas) de la
referencia, mientras que las relativas libres y semilibres introducen referentes
discursivos de naturaleza diversa (cantidad, lugar, manera, persona, tiempo). Es
así que, como plantea \textcite[327]{ditullio2014}, las consideradas
\textit{adverbiales propias} (temporales relativas, locativas, modales y de
cantidad) «denotan entidades individuales específicas como los momentos y los
lugares, así como también las maneras en que puede realizarse una cierta acción o
el grado o la intensidad con que se lleva a cabo». Dentro de este panorama, las
temporales introducidas por adverbios o preposiciones no relativas y las
tradicionalmente consideradas \textit{adverbiales impropias} ---causales,
finales, condicionales y concesivas--- suman valores que no dependen de la
estructura argumental ni del establecimiento de la referencia nominal, sino que
derivan de las relaciones lógicas de temporalidad y causalidad establecidas entre
los mundos representados en la oración principal y la subordinada.

Las construcciones con valores lógico-relacionales (temporales, causales,
finales, condicionales, concesivos) funcionan como modificadores adjuntos, ya sea
en el nivel del predicado (adjuntos circunstanciales) o en el nivel oracional
(modificadores del núcleo oracional y modificadores de la modalidad oracional).
Asimismo, establecen correlaciones temporales y modales con la oración principal
para orientar, junto con conectores y elementos conjuntivos, la interpretación
global de la oración compuesta. En este sentido, por ejemplo, la constitución de
un acto de habla compromisorio (advertencia o amenaza) puede derivar de la
interpretación condicional que surge en cualquiera de los ejemplos de
\exref{ex:13-gritando}.

\begin{exe}
  \ex\label{ex:13-gritando} Correlaciones temporales y modales más conectores y
      elementos conjuntivos:
    \begin{xlist}
      \ex\label{ex:13-gritando-a} \textit{Si me seguís gritando en la calle}, voy
          y te hago una denuncia.
      \ex\label{ex:13-gritando-b} \textit{Como me sigas gritando en la calle},
          voy y te hago una denuncia.
      \ex\label{ex:13-gritando-c} \textit{Donde me sigas gritando en la calle},
          voy y te hago la denuncia.
      \ex\label{ex:13-gritando-d} \textit{Gritándome cosas en la calle}, solo vas
          a conseguir que te haga la denuncia.
    \end{xlist}
\end{exe}

Esto es posible gracias, por un lado, a la presencia de conjunciones, ya sea la
típica de las condicionales, \forma{si} \exref{ex:13-gritando-a}, o también
otras, como \forma{como} \exref{ex:13-gritando-b} e incluso \forma{donde}
\exref{ex:13-gritando-c}, que no resultan características de prótasis
condicionales, y, por el otro, al mantenimiento de las correlaciones temporales y
modales entre las oraciones intervinientes. Estos mismos valores
lógico-relacionales se pueden obtener incluso sin encabezadores léxicos, a través
de estructuras oracionales organizadas alrededor de formas no finitas del verbo,
como en \exref{ex:13-gritando-d}.

A continuación, según las funciones sintácticas que pueden cumplir, vamos a
considerar tres subgrupos de construcciones con valores lógico-relacionales: a)
construcciones temporales que funcionan en el nivel del predicado como
circunstanciales (en este apartado exceptuamos las de \forma{cuando}, que ya
fueron incluidas con las relativas); b) causales y finales, que pueden funcionar
en el nivel del predicado como circunstanciales y en el nivel oracional como
modificadoras de la modalidad; c) condicionales y concesivas, que solo funcionan
en el nivel oracional como modificadoras del núcleo oracional o de la modalidad.

\subsubsection{Construcciones temporales}\label{sec:13-temporales}

Las construcciones que veremos en este apartado establecen una relación temporal
entre el evento manifestado en la oración principal y el presentado en la
subordinada. Manifiestan diversos valores semánticos relacionados con la conexión
básica de \sig{tiempo} \parencite{rudolph1996}, los cuales dependen del
significado del adverbio, preposición o locución temporal que las precede o de la
conjunción que las introduce. Pueden expresar simultaneidad
(\forma{mientras}\footnote{La naturaleza relativa o conjuntiva de
\forma{mientras} se encuentra en debate; no obstante, dada la índole
introductoria de este libro, no entraremos en estas cuestiones.}), anterioridad
(\forma{antes (de) que}) / posterioridad (\forma{apenas, después (de) que}),
delimitación (\forma{desde que / hasta que}), reiteración (\forma{siempre que,
cada vez que}), progresión (\forma{a medida que}), entre otros valores posibles.

Las construcciones temporales pueden cumplir funciones circunstanciales dentro
del predicado y, según sean requeridas o no por el verbo, se consideran
argumentales o no, es decir que se desempeñan como complementos o adjuntos
circunstanciales, como en \exref{ex:13-temp}.

\begin{exe}
  \ex\label{ex:13-temp} Valores temporales:
    \begin{xlist}
      \ex Podemos tomar un café \textit{\textbf{mientras} discutimos la propuesta
          laboral}. (Simultaneidad).
      \ex Arreglemos el orden del día \textit{\textbf{antes} de que empiece la
          reunión de directorio}. (Anterioridad).
      \ex Partimos para Mar del Plata \textit{\textbf{después} de que amanezca /
          \textbf{apenas} salga el sol}. (Posterioridad).
      \ex Mi hermano trabaja todo el día, \textit{\textbf{desde} que amanece
          \textbf{hasta} que anochece}. (Delimitación).
      \ex Me pongo a bailar \textit{\textbf{siempre} que escucho este tema}.
          (Reiteración).
      \ex Entendí el funcionamiento del cerebro \textit{\textbf{a medida} que
          avanzaba en la carrera}. (Progresión).
    \end{xlist}
\end{exe}

\begin{recuadro}[Una construcción particular: \textit{hacer} referido a \sig{transcurso de tiempo}]
El verbo \forma{hacer} seguido de un complemento temporal, como en \forma{Se fue
hace dos semanas}, constituye, según Martínez García (1996,
41),\pendiente{Martínez García (1996) no figura en la bibliografía} «una
estructura especial dentro de las subordinadas temporales durativas». En
principio, se trata de la única construcción del español con verbo conjugado en
indicativo que se subordina al verbo de la oración principal sin nexo:

\begin{exe}
  \ex\label{ex:13-hace1} Salieron para el aeropuerto [Ø \textit{hace dos
      horas}].
\end{exe}

Otra característica peculiar de esta construcción, que señala \sig{transcurso de
tiempo}, es que realiza un señalamiento temporal inverso, desde el presente hacia
el pasado, lo que permite explicar que \forma{hacer} en presente pueda
combinarse con verbos en pasado como \forma{salieron} \exref{ex:13-hace1} o
\forma{vivía} \exref{ex:13-hace2}:

\begin{exe}
  \ex\label{ex:13-hace2} \textit{Hace mucho tiempo} vivía en un lejano país.
\end{exe}

Además, el verbo \forma{hacer} conforma un esquema impersonal en tercera persona
singular, que presenta dos construcciones (con y sin \forma{que}) en apariencia
semejantes, aunque entre ellas existen «diferencias tanto semánticas como
sintácticas» \parencite[334]{ditullio2014}:

\begin{exe}
  \ex\label{ex:13-hace-I} \textit{Hace dos horas} salieron para el aeropuerto.
      (Tipo I: \textit{hace}\etq{temp} sin \textit{que}).
  \ex\label{ex:13-hace-II} \textit{Hace dos horas} que salieron para el
      aeropuerto. (Tipo II: \textit{hace}\etq{temp} más \textit{que}).
\end{exe}

Estas diferencias, estudiadas en \textcite{giammatteo2002}, muestran que en el
Tipo~I los constituyentes presentan un comportamiento unitario: no aceptan
permutación interna (\juicio{*}{\forma{Dos horas hace salieron para el
aeropuerto}}), pueden desplazarse conjuntamente dentro de la principal
(\forma{Salieron, hace dos horas, para el aeropuerto; Salieron para el aeropuerto
hace dos horas}), pueden usarse independientemente en respuesta a una pregunta
(\forma{¿Cuándo salieron para el aeropuerto? Hace dos horas} /
\juicio{*}{\forma{Dos horas}}),\pendiente{El original dice «no pueden usarse
independientemente», pero el ejemplo muestra lo contrario} pueden funcionar como
complemento de las preposiciones \forma{de, desde} y \forma{hasta} (\forma{Hasta
hace un rato estaban en casa}) y pueden ser conmutados por un adverbio
(\forma{Recientemente salieron para el aeropuerto}); por otra parte,
\forma{hacer} está prácticamente inmovilizado en presente. Las características
señaladas muestran que esta construcción ha sufrido un proceso de
gramaticalización que implica una descategorización: en este caso, una estructura
originalmente oracional ya no introduce otro evento diferente del de la
principal, sino que pasa a cumplir una función adverbial y a desempeñarse como un
adjunto temporal que sitúa temporalmente el suceso.

En el Tipo~II, el comportamiento de los constituyentes es diferente: pueden
permutarse internamente (\forma{Dos horas hace que salieron para el aeropuerto}),
no pueden desplazarse dentro de la principal ni tampoco omitirse
(\juicio{*}{\forma{Que salieron para el aeropuerto hace dos horas}}), pueden ser
usados independientemente en respuesta a una pregunta (\forma{¿Cuánto hace que
salieron? Dos horas}) y no funcionan como complemento de preposición ni pueden
ser reemplazados por un adverbio (\juicio{*}{\forma{Recientemente que salieron
para el aeropuerto}}); por otra parte, \forma{hacer} admite mucha más variación
en el uso de tiempos verbales.

Sintácticamente, mientras en el Tipo~I, sin \forma{que}, la construcción de
\forma{hacer} queda incluida como un adjunto en la oración principal, el Tipo~II
ejemplifica la situación inversa: la construcción de \forma{hacer}\etq{temp}
constituye la oración principal, a la que el evento de \forma{salir para el
aeropuerto} se subordina mediante la conjunción \forma{que}. Esta conformación
pone de relieve el \sig{transcurso}, mientras que el evento, expresado en la
completiva de \forma{que}, indica el momento a partir del cual se mide el
transcurso. De este modo, en \exref{ex:13-hace-II} interpretamos \sig{Transcurrieron
dos horas desde que salieron para el aeropuerto}. Por su valor semántico, esta
construcción constituye una subordinada adverbial que, como sostenía
\textcite[§\,778]{bello2002}, lleva omitida la preposición inicial \forma{de} o
\forma{desde} con el significado de \sig{a partir de} (cf. \forma{de tres a
cinco; de Victoria a Retiro}). Por lo general, la preposición se omite, aunque
puede eventualmente reponerse (\forma{Hace dos horas desde que salieron para el
aeropuerto}) y resulta obligatoria tanto cuando se pregunta por el complemento
temporal (\forma{¿De qué hace dos horas?} / \juicio{*}{\forma{¿Qué hace dos
horas?}}) como en la construcción alternativa en que la subordinada se reemplaza
por una nominalización (\forma{Hace dos horas de su salida}). Por último, solo la
construcción del Tipo~II admite otra especificación temporal inicial que indica
el límite final del transcurso (\forma{Hoy hace un año que se murió}).
\end{recuadro}

\subsubsection{Construcciones causales y finales}\label{sec:13-causales}

Las construcciones causales y finales manifiestan, en términos aristotélicos,
respectivamente, la \sig{causa eficiente} y la \sig{causa final} (o propósito).
Las causales pueden indicar la \sig{causa real} de lo que se sostiene en la
oración principal \exref{ex:13-cf-a} o una \sig{causa ilocutiva hipotética}, que
le permite al hablante hacer una inferencia respecto de lo que se dice en la
principal \exref{ex:13-mmo-a}. Las finales indican la \sig{causa final} o el
objetivo por el que se realiza un evento. Tienen carácter prospectivo, es decir
que se orientan a un tiempo posterior al de la principal, y también pueden
indicar un \sig{propósito real} \exref{ex:13-cf-b} o un \sig{propósito
ilocutivo}, que motiva al hablante a decir lo que dice en la principal
\exref{ex:13-mmo-d}.

Las oraciones subordinadas causales y finales pueden formar parte de
construcciones que cumplen funciones argumentales y no argumentales dentro del
predicado, es decir que pueden funcionar como adjuntos o complementos
circunstanciales \exref{ex:13-cf}.

\begin{exe}
  \ex\label{ex:13-cf} Circunstancial (complemento o adjunto verbal):
    \begin{xlist}
      \ex\label{ex:13-cf-a} Causa (construcción causal interna al predicado):\\
          {}[\etq{O} Habla inglés [\etq{SP/AC} por [\etq{OC/CT} que se lo
          enseñaron en la casa]]]. (Habla inglés \textit{porque se lo enseñaron
          en la casa}).
      \ex\label{ex:13-cf-b} Fin (construcción final interna al predicado):\\
          {}[\etq{O} Habla despacio [\etq{SP/AC} para [\etq{OC/CT} que lo
          entiendan bien]]]. (Habla despacio \textit{para que lo entiendan
          bien}).
      \ex\label{ex:13-cf-c} {}[\etq{O} Esta tijera sirve [\etq{SP/CC} para
          [\etq{OC/CT} cortar el césped]]]. (Esta tijera sirve \textit{para
          cortar el césped}).
    \end{xlist}
\end{exe}

En \exref{ex:13-cf-a}, la construcción causal formada por el sintagma
preposicional con núcleo en \forma{por} más la subordinada encabezada por
\forma{que} como su complemento o término es un adjunto circunstancial de causa,
es decir, un modificador del predicado que expresa la \sig{causa real} del estado
de cosas manifestado en la oración principal. En \exref{ex:13-cf-b}, la
construcción final, formada por el sintagma preposicional nucleado en
\forma{para} más la subordinada encabezada por \forma{que} como su complemento o
término, es un adjunto circunstancial de fin, un modificador del predicado que
expresa el propósito de la acción realizada. En \exref{ex:13-cf-c}, la
construcción final, formada por el sintagma preposicional nucleado en
\forma{para} más la subordinada de infinitivo como su complemento o término, es
un complemento circunstancial exigido por la semántica del verbo \forma{servir}.
Recordemos que es habitual que en las gramáticas se considere que \forma{porque}
es una conjunción que encabeza una subordinada causal y que \forma{para que} es
una locución conjuntiva que encabeza una subordinada final; sin embargo, la
posición aquí adoptada es la del enfoque analítico, que presentamos en
§\ref{sec:13-encabezadores}, dado que lo consideramos más simple y
regularizador.

Por otro lado, las construcciones con valores causales y finales también pueden
cumplir funciones no argumentales como adjuntos de nivel oracional, es decir,
como modificadores de la modalidad oracional, que constituyen funciones
periféricas dentro de la oración \exref{ex:13-mmo}.

\begin{exe}
  \ex\label{ex:13-mmo} Modificador de la modalidad oracional (adjunto
      oracional):
    \begin{xlist}
      \ex\label{ex:13-mmo-a} Causa (construcción causal externa al predicado):\\
          {}[\etq{O} Los chicos se habrán intoxicado en la cena, [\etq{SP/MMO}
          por [\etq{OC/CT} que faltaron a la escuela]]].
      \ex\label{ex:13-mmo-b} {}[\etq{O} [\etq{OS/MMO} Como tenemos trabajo
          atrasado], mañana lleguen temprano a la oficina].
      \ex\label{ex:13-mmo-c} {}[\etq{O} Mañana lleguen temprano a la oficina,
          [\etq{OS/MMO} que tenemos trabajo atrasado]].
      \ex\label{ex:13-mmo-d} Fin (construcción final externa al predicado):\\
          {}[\etq{O} [\etq{SP/MMO} Para [\etq{OC/CT} que quede constancia]], los
          estudiantes se quejan del reglamento].
      \ex\label{ex:13-mmo-e} {}[\etq{O} Ponete algo de abrigo, [\etq{SP/MMO} para
          [\etq{OC/CT} que no te enfermes]]].
    \end{xlist}
\end{exe}

La serie \exref{ex:13-mmo} muestra los valores epistémicos e ilocutivos que
pueden asumir las construcciones causales y finales cuando se orientan a la
enunciación y cambian la modalidad de la oración principal. En
\exref{ex:13-mmo-a}, la construcción causal modifica a la oración principal
expresando la razón que tiene el hablante para hacer una inferencia o conjetura,
con valor epistémico (\sig{porque faltaron a la escuela, \textit{deduzco} que los
chicos se intoxicaron en la cena}). En \exref{ex:13-mmo-b} y
\exref{ex:13-mmo-c}, la causal explicativa, antepuesta (\forma{como}) o
pospuesta (\forma{que}) a la oración principal, ofrece una explicación o una
justificación: al manifestar una atenuación se gana en cortesía (\sig{porque
tenemos trabajo atrasado, les \textit{pido} que mañana lleguen temprano a la
oficina}). De manera paralela, en \exref{ex:13-mmo-d}, la construcción final
modifica a la oración principal expresando la finalidad que motiva al hablante,
con valor ilocutivo (\sig{para que quede constancia, \textit{afirmo} que los
estudiantes se quejan del reglamento}), y en \exref{ex:13-mmo-e}, la final
modifica a la oración principal al dar una justificación que atenúa en forma
cortés el acto de habla (\sig{para que no te enfermes, \textit{sugiero} que te
pongas algo de abrigo}).

\subsubsection{Construcciones condicionales y concesivas}\label{sec:13-condicionales}

Condicionales y concesivas expresan otras dos variantes de la causa. En el caso
de las condicionales, manifiestan que la realización de lo enunciado en la
principal depende de una condición o requisito cuyo cumplimiento constituye una
\sig{causa virtual} o \sig{potencial} para que ocurra lo enunciado en la oración
principal. Las concesivas, en cambio, expresan una \sig{objeción} o dificultad
para la realización del evento manifestado en la principal que, sin embargo, no
impide su realización, por lo cual constituyen una \sig{causa inoperante} o
\sig{ineficaz}, es decir, una causa que no logra su objetivo; de modo que las
concesivas aúnan dos significados: \sig{causa} y \sig{contraste}
\parencite{giammatteo2019}. Si se confrontan las concesivas con las
condicionales, también se las puede interpretar como una condición negada, es
decir, como una «contra-causa en cuanto que el efecto es contrario al esperado»
\parencite[344]{fernandezlagunilla1995}.

Condicionales y concesivas cumplen funciones periféricas en la oración como
modificadoras del núcleo oracional o modificadoras de la modalidad oracional.
Junto con la principal forman construcciones binarias: los \termino{períodos
hipotéticos} condicional y concesivo, según se ejemplifica en
\exref{ex:13-mno}, \exref{ex:13-mmocond} y \exref{ex:13-mmoconc}.

\begin{exe}
  \ex\label{ex:13-mno} Modificador del núcleo oracional (adjunto oracional):
    \begin{xlist}
      \ex\label{ex:13-mno-a} Condición (período hipotético condicional):\\
          {}[\etq{O} [\etq{OS/MNO} Si te ofrecen algo de beber], no vayas a
          aceptarlo].
      \ex\label{ex:13-mno-b} Concesión (período hipotético concesivo):\\
          {}[\etq{O} [\etq{SAdv/MNO} Aun [\etq{OC/C} que te ofrezcan algo de
          comer]], no lo aceptes]. (\textit{Aunque te ofrezcan algo de comer},
          no lo aceptes).
    \end{xlist}
\end{exe}

Las oraciones de \exref{ex:13-mno} presentan los períodos hipotéticos que se
constituyen entre las construcciones condicionales \exref{ex:13-mno-a} o
concesivas \exref{ex:13-mno-b} y la oración principal. En \exref{ex:13-mno-a},
la \termino{prótasis} condicional encabezada por \forma{si} manifiesta una causa
hipotética del evento contenido en la \termino{apódosis}, que condiciona su
realización. En \exref{ex:13-mno-b}, la prótasis concesiva formada por el
sintagma adverbial nucleado por el adverbio \forma{aun} más la subordinada
encabezada por \forma{que} como su complemento manifiesta una \sig{causa
ineficaz} del evento contenido en la apódosis, que resulta contrario a las
expectativas. De modo semejante a lo que sucede con \forma{porque}, en las
gramáticas se suele considerar que \forma{aunque} es una conjunción que encabeza
una subordinada concesiva; no obstante, en la versión analítica que aquí
presentamos, la subordinada introducida por \forma{que} es una completiva que
funciona como complemento del adverbio \forma{aun}.

\begin{exe}
  \ex\label{ex:13-mmocond} Modificador de la modalidad oracional (adjunto
      oracional). Condición (período hipotético condicional):
    \begin{xlist}
      \ex\label{ex:13-mmocond-a} {}[\etq{O} [\etq{OS/MMO} Si adornó toda la casa],
          hoy es su cumpleaños y recibe a sus amigas].
      \ex\label{ex:13-mmocond-b} {}[\etq{O} [\etq{OS/MMO} Si no me falla la
          memoria], hoy es el cumpleaños de tu hermano].
      \ex\label{ex:13-mmocond-c} {}[\etq{O} [\etq{OS/MMO} Si te parece bien],
          celebremos su cumpleaños aunque falte un día].
      \ex\label{ex:13-mmocond-d} {}[\etq{O} [\etq{OS/MMO} Si consideramos la
          fecha], hoy deberíamos celebrar su cumpleaños].
    \end{xlist}
\end{exe}

Las oraciones de \exref{ex:13-mmocond} muestran los valores epistémicos e
ilocutivos de los períodos condicionales que se orientan a la enunciación: la
prótasis condicional modifica la modalidad de la oración principal. En
\exref{ex:13-mmocond-a}, la prótasis manifiesta una condición para inferir una
conjetura o una conclusión (el evento contenido en la apódosis) y el período
condicional asume un valor epistémico inferencial (\sig{si adornó toda la casa,
\textit{deduzco} que hoy es su cumpleaños y recibe a sus amigas}). En
\exref{ex:13-mmocond-b}, al igual que en \exref{ex:13-mmocond-c} y
\exref{ex:13-mmocond-d}, la prótasis modifica la fuerza ilocutiva (el acto de
habla que forma la apódosis) y el período condicional asume valores ilocutivos:
manifiesta una atenuación (\sig{si no me falla la memoria, \textit{afirmo} que
hoy es el cumpleaños de tu hermano}), imprime cortesía (\sig{si te parece bien,
\textit{propongo} que celebremos su cumpleaños, aunque falte un día}) o bien
orienta sobre el acto lingüístico (\sig{si consideramos la fecha, \textit{digo}
que hoy deberíamos celebrar su cumpleaños}).

\begin{exe}
  \ex\label{ex:13-mmoconc} Modificador de la modalidad oracional (adjunto
      oracional). Concesión (período hipotético concesivo):
    \begin{xlist}
      \ex\label{ex:13-mmoconc-a} {}[\etq{O} [\etq{SAdv/MMO} Aun [\etq{OC/C} que
          haya pruebas en su contra]], legalmente debe ser considerado
          inocente].
      \ex\label{ex:13-mmoconc-b} {}[\etq{O} [\etq{SAdv/MMO} Aun [\etq{OC/C} que
          parezca mentira]], es la primera vez que lo cita la justicia].
      \ex\label{ex:13-mmoconc-c} {}[\etq{O} [\etq{SAdv/MMO} Aun [\etq{OC/C} que
          no entiendo cómo sucedió]], tu hermano no tuvo la culpa].
    \end{xlist}
\end{exe}

Las oraciones de \exref{ex:13-mmoconc} muestran los valores epistémicos e
ilocutivos de los períodos concesivos orientados a la enunciación: la prótasis
concesiva modifica la modalidad oracional de la oración principal. En
\exref{ex:13-mmoconc-a}, la prótasis manifiesta una condición ineficaz para
inferir determinada conjetura o conclusión respecto del evento contenido en la
apódosis y el período concesivo toma un valor epistémico (\sig{aunque haya
pruebas en su contra, \textit{supongo} que legalmente es inocente}). En
\exref{ex:13-mmoconc-b} y \exref{ex:13-mmoconc-c}, la prótasis modifica la fuerza
ilocutiva (el acto de habla de la apódosis) y el período concesivo toma un valor
ilocutivo: atenúa el acto lingüístico (\sig{aunque parezca mentira,
\textit{afirmo} que es la primera vez que lo cita la justicia}) y ofrece una
orientación (\sig{aunque no entiendo cómo sucedió, \textit{digo} que tu hermano
no tuvo la culpa}).

\begin{recuadro}[Las construcciones con valores adverbiales con formas verbales no finitas]
Según lo señalado en §\ref{sec:13-encabezador}, algunas oraciones subordinadas se
construyen con las formas no finitas del verbo (infinitivo, participio y
gerundio). Estas subordinadas nunca van introducidas por la conjunción
\forma{que}, que requiere un verbo conjugado en la subordinada. Las de participio
y gerundio no llevan encabezadores obligatorios, aunque pueden admitir alguno
opcional: \forma{\textbf{Apenas} terminada la reunión, nos dirigimos al
aeropuerto; \textbf{Aun} considerando el mejor escenario, el rinde de la cosecha
va a resultar pobre}. Las de infinitivo pueden construirse sin encabezador o
formar parte de una construcción preposicional, en la que la preposición no
puede omitirse porque explicita el valor semántico de la subordinada:
\forma{El cielo se tiñó de rojo \textbf{al} caer la tarde; Recibió media falta
\textbf{por} haber llegado tarde / llegar tarde; Nos acercamos a Secretaría
\textbf{para} retirar el diploma}.

Una característica de las subordinadas con verbo no flexionado es que el sujeto
puede coincidir con el sujeto de la principal (como en \forma{Teniendo tiempo,
voy a revisar el trabajo}) o no coincidir con este (como en \forma{Muerto el
perro, se acabó la rabia}). Cuando el sujeto de la subordinada difiere del de la
principal, se pospone al verbo y la subordinada se suele separar de la oración
principal por una pausa (representada por coma en la escritura). Debido a estas
características, que les otorgan mayor independencia respecto de la principal,
estas subordinadas se han denominado \termino{absolutas} (desligadas):
\forma{Agotados los temas del orden del día, se da por concluida la asamblea; Se
da por concluida la asamblea una vez agotados los temas del orden del día}.

Según muestran los ejemplos anteriores, estas subordinadas pueden manifestar
diversos valores adverbiales. Cuando forman parte de una construcción
preposicional (\forma{\textbf{por} haber llegado tarde, \textbf{para} retirar el
diploma}), es la preposición la que marca el valor semántico de la subordinada;
pero muchas veces, cuando se construyen como absolutas, como sostiene
\textcite[672]{morenocabrera1991}, las subordinadas con formas no finitas se
utilizan «cuando no se quiere especificar el tipo de relación que existe entre el
estado de hechos denotado por la subordinada y el denotado por la oración». Es
decir que, al no tener un elemento previo que fije la interpretación, es el
contexto el que determina en cada caso el valor semántico específico de estas
oraciones, que muchas veces puede quedar ambiguo: \forma{Muerto el perro, se
acabó la rabia} (tiempo, causa, condición).
\end{recuadro}

\subsection{Construcciones con valores cuantificacionales-intensificativos}\label{sec:13-cuantificativas}

Se trata de construcciones que funcionan como complemento de un elemento
cuantificativo (\forma{tanto, más, menos}) que modifica el grado de una categoría
léxica (adjetivo, adverbio, verbo o nombre). Forman estructuras correlativas del
tipo \forma{tanto\ldots\ como, más\ldots\ que, tan(to)\ldots\ que}, etc. Las
oraciones subordinadas comparativas y consecutivas modifican el grado de un
elemento cuantificador o intensificador correlativo.

\subsubsection{Construcciones comparativas}\label{sec:13-comparativas}

De las cuatro relaciones básicas señaladas por \textcite{rudolph1996}, las
construcciones comparativas están vinculadas con las de adición y con las de
contraste, que permiten establecer semejanzas y diferencias entre entidades,
propiedades o sucesos. Así, en una estructura comparativa se puede realizar una
comparación sobre el número de elementos de un conjunto (\forma{Tengo
\textbf{menos} exámenes en el escritorio \textbf{que alumnos en la clase}}),
sobre la cantidad de alguna sustancia (\forma{Vi \textbf{más} basura en la playa
\textbf{que el verano pasado}}) o sobre el grado de una propiedad (\forma{La
ópera es \textbf{tan} intensa \textbf{como esperaba}}). En estas construcciones
se presentan valores cuantificables o graduables sujetos a comparación y la
estructura comparativa puede ser tanto una oración subordinada como un fragmento
oracional, siempre con las conjunciones \forma{que} y \forma{como} en su
encabezado. Los ejemplos de \exref{ex:13-compar} presentan las construcciones
comparativas como estructuras cuantificativas complejas.

\begin{exe}
  \ex\label{ex:13-compar} Coda comparativa (adjunto del cuantificador):
    \begin{xlist}
      \ex\label{ex:13-compar-a} {}[\etq{O} Este año la empresa va a trabajar
          \textit{más} horas [\etq{SC/CCM} que el año pasado]].
      \ex\label{ex:13-compar-b} {}[\etq{O} Este año la empresa tiene \textit{mucho
          menos} trabajo [\etq{SC/CCM} que el año pasado]].
      \ex\label{ex:13-compar-c} {}[\etq{O} Este año la empresa desarrolló un
          trabajo \textit{más} intenso [\etq{SC/CCM} que el año pasado]].
      \ex\label{ex:13-compar-d} {}[\etq{O} Este año la empresa trabajó
          \textit{tan} intensamente [\etq{SC/CCM} como el año pasado]].
      \ex\label{ex:13-compar-e} {}[\etq{O} Este año la empresa produjo \textit{dos
          veces más} [\etq{OC/CCM} de lo que se proyectó]].
    \end{xlist}
\end{exe}

\pendiente{Siglas CCM (coda comparativa) y CCN (coda consecutiva): se solapan con
otras abreviaturas; consultar a los autores}

Las oraciones de \exref{ex:13-compar} muestran la comparación del número en los
nombres contables (\forma{horas} en \exref{ex:13-compar-a}) y la comparación de
la cantidad en los nombres incontables (\forma{trabajo} en
\exref{ex:13-compar-b}). Asimismo, estas construcciones también permiten la
comparación del grado en los atributos y en los predicados, señalando el grado
del adjetivo \forma{intenso} en \exref{ex:13-compar-c}, el grado del adverbio
\forma{intensamente} en \exref{ex:13-compar-d} y el grado del verbo
\forma{producir} en \exref{ex:13-compar-e}. Tomados en conjunto, estos ejemplos
muestran que el resultado de la comparación puede indicar una desigualdad
(contraste), si los valores sujetos a comparación no coinciden
(\forma{más / menos\ldots\ que / de}), o una igualdad (semejanza), si los valores
sujetos a comparación coinciden (\forma{tan / tanto\ldots\ como}).

\subsubsection{Construcciones consecutivas}\label{sec:13-consecutivas}

Las construcciones consecutivas están vinculadas con la noción de causalidad.
Según \textcite[343]{ditullio2014}, «desde un punto de vista que solo contemple
las relaciones semánticas de orden lógico, estas oraciones son las inversas de
las causales, ya que expresan el resultado o efecto, de lo expresado por la
principal, la causa». Así, Di Tullio compara dos oraciones, \forma{Estoy agotado
porque caminé mucho} y \forma{Caminé tanto que estoy agotado}, para concluir que
difieren por el carácter ponderativo de la segunda, ausente en la primera, que
deriva de la presencia del intensificador \forma{tanto}.

En una estructura consecutiva se realiza una intensificación cuantitativa
(\forma{Había \textbf{tanta} basura en la playa \textbf{que no se podía
caminar}}) o cualitativa (\forma{Teníamos \textbf{tal} gentío en la puerta
\textbf{que debimos entregar números}}) y la consecuencia de esta intensificación
es un evento, manifestado por una subordinada encabezada por la conjunción
subordinante \forma{que}. Los ejemplos de \exref{ex:13-consec}, por su parte,
muestran las construcciones consecutivas como estructuras cuantificativas
complejas.

\begin{exe}
  \ex\label{ex:13-consec} Coda consecutiva (adjunto del cuantificador):
    \begin{xlist}
      \ex\label{ex:13-consec-a} {}[\etq{O} Tuvimos \textit{tantos} despachos
          [\etq{OC/CCN} que contratamos más empleados]].
      \ex\label{ex:13-consec-b} {}[\etq{O} Tuvimos \textit{tanta} inactividad
          [\etq{OC/CCN} que casi cerramos el local]].
      \ex\label{ex:13-consec-c} {}[\etq{O} El trabajo estuvo \textit{tan} flojo
          [\etq{OC/CCN} que redujimos los turnos a la mitad]].
      \ex\label{ex:13-consec-d} {}[\etq{O} Trabajé \textit{tan} bien [\etq{OC/CCN}
          que todavía no lo puedo creer]].
      \ex\label{ex:13-consec-e} {}[\etq{O} La empresa facturó \textit{tanto}
          [\etq{OC/CCN} que en un mes se alcanzó la meta bimestral]].
      \ex\label{ex:13-consec-f} {}[\etq{O} Las ventas fueron de \textit{tal /
          tanta} magnitud [\etq{OC/CCN} que reprogramamos muchos envíos]].
    \end{xlist}
\end{exe}

La serie \exref{ex:13-consec} presenta la intensificación cuantitativa del número
en los nombres contables (\forma{despachos} en \exref{ex:13-consec-a}) y de la
cantidad en los nombres incontables (\forma{inactividad} en
\exref{ex:13-consec-b}). Asimismo, estas construcciones permiten la
intensificación del grado en los atributos, en los predicados y en nombres
abstractos: el grado del adjetivo \forma{flojo} en \exref{ex:13-consec-c}, el
grado del adverbio \forma{bien} en \exref{ex:13-consec-d}, el grado del verbo
\forma{facturar} en \exref{ex:13-consec-e} y el grado del nombre \forma{magnitud}
en \exref{ex:13-consec-f}. En todos los casos, los ejemplos presentan la
intensificación extrema de lo expresado por una palabra léxica (nombre,
adjetivo, adverbio, verbo) como causa de la consecuencia codificada en la oración
subordinada consecutiva (\forma{tal / tan / tanto\ldots\ que}).

En las construcciones comparativas se da la correlación entre un cuantificador y
la coda comparativa y, paralelamente, en el caso de las consecutivas, la
correlación se establece entre un intensificador y la coda consecutiva. En cuanto
a las estructuras internas de tales codas, solo la coda comparativa admite tanto
oraciones subordinadas como sintagmas no oracionales, mientras que la coda
consecutiva solo admite oraciones subordinadas.

\begin{ejercicio}\label{ej:13-paratextos}
En los siguientes textos breves, formados por secuencias paratextuales (volanta,
título, copete) de noticias periodísticas, identifique las oraciones
subordinadas y las funciones sintácticas que cumplen, según se muestra en el
ejemplo resuelto~1.

% Textos periodísticos protegidos: se incorporan a mano desde el original.
\textoaincorporar{los diez paratextos periodísticos (volanta, título y copete)
del original: 1. Guido Kaczka y «Los 8 escalones» (\textit{La Razón},
21/07/2015); 2. Cacho Castaña en el Colón (\textit{La Razón}, 05/12/2016); 3.
camiones para tapar baches (\textit{La Razón}, 01/06/2016); 4. carta de Ángel Di
María (\textit{Infobae}, 05/08/2019); 5. declaraciones de Diego Maradona
(\textit{La Razón}, 22/07/2015); 6. la Selección en Houston (\textit{La Razón},
22/06/2016); 7. Jazmín Stuart (\textit{La Razón}, 24/11/2016); 8. «Maradona sin
vueltas» (\textit{La Razón}, 22/07/2015); 9. feriado por Güemes (\textit{La
Razón}, 01/06/2016); 10. fiesta clandestina en Moreno (\textit{La Razón},
22/09/2016)}

\begin{resuelto}
1. La OR encabezada por el pronombre relativo \forma{que} (\forma{que no le
gustó nada a Guido Kaczka}) es un adjunto especificativo (AS) del antecedente
nominal \forma{broma}. La OC de infinitivo (sin subordinante, \forma{participar})
es el complemento o término (CT) de la preposición \forma{a} (valor final). La
OC encabezada por la conjunción \forma{que} (\forma{que se cansó y se fue}) es la
coda consecutiva (CCN) correlativa del intensificador \forma{tanto}, con el que
forma una estructura cuantificativa compleja.
\end{resuelto}
\end{ejercicio}

\begin{recuadro}[Reflexión final]
Entre las relaciones interoracionales, la subordinación (hipotaxis) permite crear
sentidos complejos al vincular mundos completos (reales, irreales, posibles,
pasados, eventuales) que, en principio, no presentan vínculos entre sí. Más
estructurales o más periféricas, las relaciones de subordinación caracterizan a
la oración compuesta de la misma manera en que las relaciones de concordancia
caracterizan a la oración simple. Miradas en conjunto, las oraciones subordinadas
siguen patrones estructurales notablemente regulares: encabezadores léxicos
conjuntivos o relativos con verbos finitos, por un lado; verbos no finitos sin
encabezadores léxicos conjuntivos o relativos, por el otro. Los valores
semántico-pragmáticos que estas oraciones adoptan, al mismo tiempo, se pueden
resumir en dos ejes: el de la temporalidad (adición / tiempo) y el del punto de
vista (contraste / causalidad). Con estos pocos recursos gramaticales
articulamos la subjetividad de nuestros pensamientos complejos y, si acaso, los
convertimos en conocimientos intersubjetivos cuando los comunicamos a personas
que, aun sin conocerlas, tenemos por iguales.
\end{recuadro}

% --- Secciones finales -------------------------------------------------
\lecturasrecomendadas

\solucionario

\begin{solucion}{\ref{ej:13-paremias}}
3. [\etq{Ocomp} [\etq{OS} \textit{Cuando el invierno primavea}], la primavera
invernea]. 4. [\etq{Ocomp} [\etq{OS} \textit{Cuando veas} [\etq{OS} \textit{las
barbas de tu vecino cortar}]], pon las tuyas a [\etq{OS} \textit{remojar}]]. 5.
[\etq{Ocomp} [\etq{OS} \textit{Dando gracias por agravios}] negocian los hombres
sabios]. 6. [\etq{Ocomp} [\etq{OS} \textit{Diciendo las verdades}] se pierden
amistades]. 7. [\etq{Ocomp} [\etq{OS} \textit{Donde fueres}], haz [\etq{OS}
\textit{lo que vieres}]]. 8. [\etq{Ocomp} [\etq{OS} \textit{Donde hubo fuego}],
cenizas quedan]. 9. [\etq{Ocomp} El [\etq{OS} \textit{salir de la posada}] es la
mejor jornada]. 10. [\etq{Ocomp} Es de bien nacidos [\etq{OS} \textit{ser
agradecidos}]]. 11. [\etq{Ocomp} [\etq{OS} \textit{Hablando del rey de Roma}], por
la puerta asoma]. 12. [\etq{Ocomp} [\etq{OS} \textit{Lo que natura non da}]
Salamanca non presta]. 13. [\etq{Ocomp} No por [\etq{OS} \textit{mucho
madrugar}] amanece más temprano]. 14. [\etq{Ocomp} Nunca es tarde [\etq{OS}
\textit{si la dicha es buena}]]. 15. [\etq{Ocomp} Pájaro [\etq{OS} \textit{que
comió}], voló]. 16. [\etq{Ocomp} Perro [\etq{OS} \textit{que ladra}] no muerde].
17. [\etq{Ocomp} [\etq{OS} \textit{Preguntando}] se llega a Roma]. 18.
[\etq{Ocomp} [\etq{OS} \textit{Si quieres} [\etq{OS} \textit{ser bien
servido}]], sírvete a ti mismo].\pendiente{El ejercicio 1 del original no
incluye el refrán 9 (\textit{El salir de la posada}), que sí aparece en los
ejercicios 2 y 6: se agregó y se renumeró}
\end{solucion}

\begin{solucion}{\ref{ej:13-introductor}}
3. Con relativo \forma{cuando}. 4. Con relativo \forma{cuando}, y sin
subordinante en [\etq{OS} \textit{las barbas de tu vecino cortar}] y [\etq{OS}
\textit{remojar}]. 5. Sin subordinante. 6. Sin subordinante. 7. Con relativos
\forma{donde} y \forma{lo que}. 8. Con relativo \forma{donde}. 9. Sin
subordinante. 10. Sin subordinante. 11. Sin subordinante. 12. Con relativo
\forma{lo que}. 13. Sin subordinante. 14. Con conjunción \forma{si}. 15. Con
relativo \forma{que}. 16. Con relativo \forma{que}. 17. Sin subordinante. 18.
Con conjunción \forma{si}, y sin subordinante en [\etq{OS} \textit{ser bien
servido}].
\end{solucion}

\begin{solucion}{\ref{ej:13-principito}}
Respuestas posibles: 2. [\etq{Ocomp} [\etq{OS} \textbf{Si} queremos disfrutar de
las mariposas], a veces tenemos que aguantar a las orugas] / [\etq{Ocomp}
[\etq{OR} \textbf{Cuando} queremos disfrutar de las mariposas], tenemos que
aguantar a las orugas]. 3. [\etq{Ocomp} Nunca te olvides de [\etq{OC}
\textbf{que} eres el dueño de tu vida y tus emociones]]. 4. [\etq{Ocomp}
[\etq{OS} \textbf{Como} los hombres ya no tienen tiempo], compran cosas hechas a
los comerciantes] / [\etq{Ocomp} Los hombres compran cosas hechas a los
comerciantes por [\etq{OC} \textbf{que} ya no tienen tiempo]]. 5. [\etq{Ocomp}
Me pregunto [\etq{OC} \textbf{si} las estrellas se iluminan con algún fin]]. 6.
[\etq{Ocomp} Nunca estamos contentos en el lugar [\etq{OR} \textbf{donde}
estamos]] / [\etq{Ocomp} Nunca estamos contentos [\etq{OR} \textbf{cuando}
estamos en un lugar]]. 7. [\etq{Ocomp} Todas las estrellas ríen [\etq{OR}
\textbf{cuando} brilla la luna]]. 8. [\etq{Ocomp} Aun [\etq{OC} \textbf{que}
pocas lo recuerdan], todas las personas grandes han sido niños antes] /
[\etq{Ocomp} Aun [\etq{OR} \textbf{cuando} pocas lo recuerdan], todas las
personas grandes han sido niños antes]. 9. [\etq{Ocomp} [\etq{OS} \textbf{Si} un
hombre es vanidoso], todos los demás hombres son admiradores] / [\etq{Ocomp}
Todos los demás hombres son admiradores [\etq{OR} \textbf{cuando} un hombre es
vanidoso]]. 10. [\etq{Ocomp} Hay que hacer cuidadosamente la limpieza del
planeta [\etq{OR} \textbf{cuando} uno termina de arreglarse]] / [\etq{Ocomp} Hay
que hacer cuidadosamente la limpieza del planeta al [\etq{OC} terminar de
arreglarse]].
\end{solucion}

\begin{solucion}{\ref{ej:13-analitica}}
2. [\etq{Ocomp} [\etq{SAdv} \textit{Aun} [\etq{OC} \textbf{que} llueva]], vamos a
salir], con valor concesivo en \forma{aun}. 3. [\etq{Ocomp} Habló con el rector
[\etq{SP} \textit{con el fin de} [\etq{OC} \textbf{que} se aprobase su curso]]],
con valor final en \forma{con el fin de} (locución preposicional). 4.
[\etq{Ocomp} Lo explico de nuevo [\etq{SP} \textit{para} [\etq{OC} \textbf{que}
no queden dudas]]], con valor final en \forma{para}. 5. [\etq{Ocomp} Me alegro
[\etq{SP} \textit{por} [\etq{OC} \textbf{que} lo hice con cariño]]], con valor
causal en \forma{por}. 6. [\etq{Ocomp} No pude llegar antes [\etq{SP}
\textit{por} [\etq{OC} \textbf{que} había mucho tráfico]]], con valor causal en
\forma{por}. 7. [\etq{Ocomp} No va a aceptar, [\etq{SAdv} \textit{aun} [\etq{OC}
\textbf{que} se lo pidamos amablemente]]], con valor concesivo en \forma{aun}. 8.
[\etq{Ocomp} Nos alegra tu visita [\etq{SP} \textit{por} [\etq{OC} \textbf{que}
te extrañábamos]]], con valor causal en \forma{por}. 9. [\etq{Ocomp} Te lo
repito [\etq{SP} \textit{para} [\etq{OC} \textbf{que} lo entiendas]]], con valor
final en \forma{para}. 10. [\etq{Ocomp} Trabajamos [\etq{SP} \textit{con el fin
de} [\etq{OC} \textbf{que} el proyecto sea un éxito]]], con valor final en
\forma{con el fin de} (locución preposicional).
\end{solucion}

\begin{solucion}{\ref{ej:13-perifrasis}}
3. [\etq{Ocomp} No voy a la fiesta [\etq{SP} por [\etq{OC} \textbf{que} estoy
cansado]]]: oración subordinada con verbo en forma personal introducida por la
conjunción subordinante \forma{que}. 4. [\etq{Ocomp} [\etq{SAdv} Aun [\etq{OS}
\textit{estando cansados}]], trabajamos duro]: oración subordinada con verbo en
forma no personal, gerundio. 5. [\etq{Ocomp} [\etq{SAdv} Aun [\etq{OC}
\textbf{que} estábamos cansados]], trabajamos duro]: oración subordinada con
verbo en forma personal introducida por la conjunción subordinante \forma{que}.
6. [\etq{Ocomp} [\etq{OS} \textit{No estando invitado}], no \textit{puedo ir} al
cumpleaños]: oración subordinada con verbo en forma no personal, gerundio; en la
principal, perífrasis verbal de \forma{poder} + infinitivo, de valor modal
deóntico. [\etq{Ocomp} [\etq{OS} \textbf{Si} no me invitan], no puedo ir al
cumpleaños]: oración subordinada con verbo en forma personal introducida por la
conjunción subordinante \forma{si}; en la principal, perífrasis verbal de
\forma{poder} + infinitivo, de valor modal deóntico. 7. [\etq{Ocomp} Es
importante [\etq{OC} \textit{comprender el funcionamiento del motor}]]: oración
subordinada con verbo en forma no personal, infinitivo. 8. [\etq{Ocomp} Es
importante [\etq{OC} \textbf{que} comprendas el funcionamiento del motor]]:
oración subordinada con verbo en forma personal introducida por la conjunción
subordinante \forma{que}. 9. [\etq{Ocomp} Logró [\etq{OC} \textbf{que} sus
amigos supieran la verdad]]: oración subordinada con verbo en forma personal
introducida por la conjunción subordinante \forma{que}. 10. [\etq{Ocomp} Logró
[\etq{OC} \textit{saber la verdad}]]: oración subordinada con verbo en forma no
personal, infinitivo. 11. [\etq{Ocomp} Quiso [\etq{OC} \textbf{que} sus amigos
supieran la verdad]]: oración subordinada con verbo en forma personal
introducida por la conjunción subordinante \forma{que}. 12. [\etq{Ocomp} Quiso
[\etq{OC} \textit{saber la verdad}]]: oración subordinada con verbo en forma no
personal, infinitivo. 13. [\etq{Ocomp} Se negó [\etq{SP} a [\etq{OC}
\textbf{que} sus amigos supieran la verdad]]]: oración subordinada con verbo en
forma personal introducida por la conjunción subordinante \forma{que}. 14.
[\etq{Ocomp} Se negó [\etq{SP} a [\etq{OC} \textit{saber la verdad}]]]: oración
subordinada con verbo en forma no personal, infinitivo. 15. [\etq{O} De esta
manera \textit{comenzó a saber} la verdad]: oración simple con perífrasis verbal
de \forma{comenzar} + \forma{a} + infinitivo, de valor aspectual incoativo. 16.
[\etq{O} De esta manera \textit{pudo saber} la verdad]: oración simple con
perífrasis verbal de \forma{poder} + infinitivo, de valor modal, capacidad. 17.
[\etq{O} De esta manera \textit{va a saber} la verdad]: oración simple con
perífrasis verbal de \forma{ir} + \forma{a} + infinitivo, de valor temporal
futuro. 19. [\etq{Ocomp} El tenor ingresó a la escena [\etq{OS} \textbf{mientras}
cantaba el aria]]: oración subordinada con verbo en forma personal introducida
por la conjunción subordinante \forma{mientras}. 20. [\etq{Ocomp} El tenor
ingresó [\etq{OS} \textit{cantando el aria}]]: oración subordinada con verbo en
forma no personal, gerundio.
\end{solucion}

\begin{solucion}{\ref{ej:13-funciones}}
3. [\etq{Ocomp} [\etq{OR/AC} \textit{Cuando el invierno primavea}], la primavera
invernea]: la OR libre encabezada por el adverbio relativo \forma{cuando} es un
adjunto circunstancial (AC) de tiempo; atendiendo al sentido generalizador del
refrán, admite otra interpretación, en la que la oración subordinada funciona
como un adjunto modificador de núcleo oracional: la prótasis de un período
hipotético condicional (\sig{si el invierno primavea, la primavera invernea}).

4. [\etq{Ocomp} [\etq{OR/AC} Cuando veas [\etq{OC/COD} \textit{las barbas de tu
vecino cortar}]], pon las tuyas a [\etq{OC/CT} \textit{remojar}]]: la OR libre
encabezada por el relativo \forma{cuando} es un adjunto circunstancial (AC) de
tiempo, y admite otra interpretación como adjunto modificador de núcleo
oracional con valor condicional; la primera OC de infinitivo (sin subordinante)
es el complemento objeto directo (COD) del verbo transitivo \forma{ver}
(\sig{cuando veas que cortan las barbas de tu vecino}) y la segunda OC de
infinitivo (sin subordinante) es el complemento o término (CT) de la preposición
\forma{a} (valor final).

5. [\etq{Ocomp} [\etq{OS/AC} \textit{Dando gracias por agravios}] negocian los
hombres sabios]: la OS de gerundio (sin subordinante) es un adjunto
circunstancial (AC) de modo o manera.

6. [\etq{Ocomp} [\etq{OS/AC} \textit{Diciendo las verdades}] se pierden
amistades]: la OS de gerundio (sin subordinante) es un adjunto circunstancial
(AC) de modo o manera.

7. [\etq{Ocomp} [\etq{OR/AC} \textit{Donde fueres}], haz [\etq{OR/COD}
\textit{lo que vieres}]]: la OR libre encabezada por el adverbio relativo
\forma{donde} es un adjunto circunstancial (AC) de lugar y la OR semilibre
encabezada por el relativo \forma{lo que} es el COD del verbo transitivo
\forma{hacer}; por su valor generalizador como refrán, admite otra
interpretación, en la que la OS encabezada por \forma{donde} con verbo en
subjuntivo (futuro) es un adjunto modificador de núcleo oracional con valor
condicional (\sig{si estuvieres afuera, haz lo que vieres}).

8. [\etq{Ocomp} [\etq{OR/AC} \textit{Donde hubo fuego}], cenizas quedan]: la OR
libre encabezada por el adverbio relativo \forma{donde} es un adjunto
circunstancial (AC) de lugar; por su valor generalizador como refrán, admite otra
interpretación, en la que la OS encabezada por \forma{donde} es un adjunto
modificador de núcleo oracional con valor condicional (\sig{si hubo fuego,
cenizas quedan}).

9. [\etq{Ocomp} El [\etq{OC/S} \textit{salir de la posada}] es la mejor
jornada]: la OC de infinitivo (sin subordinante) es el sujeto (S) de la
oración; el determinante definido aporta valor expresivo (\sig{el salir de la
posada es la mejor jornada} = \sig{salir de la posada es la mejor jornada}).

10. [\etq{Ocomp} Es de bien nacidos [\etq{OC/S} \textit{ser agradecidos}]]: la
OC de infinitivo (sin subordinante) es el sujeto (S) de la oración; el
determinante definido aportaría valor expresivo (\sig{es de bien nacidos ser
agradecidos} = \sig{es de bien nacidos el ser agradecidos}).

11. [\etq{Ocomp} [\etq{OS/MNO} \textit{Hablando del rey de Roma}], por la puerta
asoma]: la OS de gerundio (sin subordinante) es un adjunto de nivel oracional, un
modificador de núcleo oracional (MNO) que introduce un referente discursivo que
da las condiciones de verdad de toda la oración; la OS puede tener una
interpretación temporal (\sig{cuando se habla del rey de Roma, por la puerta
asoma}).

12. [\etq{Ocomp} [\etq{OR/COD} \textit{Lo que natura non da}] Salamanca non
presta]: la OR semilibre encabezada por el relativo \forma{lo que} es el
complemento objeto directo (COD) del verbo transitivo \forma{prestar}; la
interpretación global del refrán, generalizadora, es la de un período hipotético
condicional (\sig{si natura non lo da, Salamanca non lo presta}).

13. [\etq{Ocomp} No por [\etq{OC/CT} \textit{mucho madrugar}] amanece más
temprano]: la OC de infinitivo (sin subordinante) es el complemento o término
(CT) de la preposición \forma{por} (valor causal); la construcción tiene el
mismo valor contrastivo que un período hipotético concesivo (\sig{aunque mucho
se madrugue, no amanece más temprano}).

14. [\etq{Ocomp} Nunca es tarde [\etq{OS/MNO} \textit{si la dicha es buena}]]:
la OS encabezada por la conjunción subordinante \forma{si} es un adjunto de nivel
oracional, un modificador de núcleo oracional (MNO) que introduce un estado de
cosas en el mundo que da las condiciones de verdad para el período hipotético
condicional; existe una lectura temporal de fondo, sobre la que se monta la
interpretación condicional, como se aprecia en la versión alternativa del refrán
(\forma{Nunca es tarde cuando la dicha es buena}).

15. [\etq{Ocomp} Pájaro [\etq{OR/AS} \textit{que comió}], voló]: la OR
encabezada por el relativo \forma{que} es un adjunto especificativo (AS) del
antecedente nominal \forma{pájaro}, cuya referencia se restringe.

16. [\etq{Ocomp} Perro [\etq{OR/AS} \textit{que ladra}] no muerde]: la OR
encabezada por el relativo \forma{que} es un adjunto especificativo (AS) del
antecedente nominal \forma{perro}, cuya referencia se restringe.

17. [\etq{Ocomp} [\etq{OS/AC} \textit{Preguntando}] se llega a Roma]: la OS de
gerundio (sin subordinante) es un adjunto circunstancial (AC) de modo o manera;
la OS de gerundio, carente de una marca (conjunción, adverbio o preposición) que
oriente en la interpretación, puede dar lugar a lecturas diversas:
\textit{temporal} (\sig{cuando se pregunta}), \textit{condicional} (\sig{si se
pregunta}) e incluso \textit{causal} (\sig{porque se pregunta}).

18. [\etq{Ocomp} [\etq{OS/MNO} \textit{Si quieres} [\etq{OC/COD} \textit{ser bien
servido}]], sírvete a ti mismo]: la OS encabezada por la conjunción subordinante
\forma{si} es un adjunto de nivel oracional, un modificador de núcleo oracional
(MNO) que introduce un estado de cosas en el mundo que da las condiciones de
verdad para el período hipotético condicional; la OC de infinitivo (sin
subordinante) es el complemento objeto directo (COD) del verbo transitivo
\forma{querer}.
\end{solucion}

\begin{solucion}{\ref{ej:13-relativos}}
2. [\etq{O} Compré una valija]. [\etq{Ocomp} De hecho, compré [\etq{OR/COD}
\textit{la que} me recomendaste]]: \forma{la que}, complemento objeto directo del
verbo transitivo \forma{recomendar}; antecedente contextual, \forma{valija}. 4.
[\etq{Ocomp} Decime [\etq{OR/COD} \textit{lo que} pensás]]: \forma{lo que},
complemento objeto directo del verbo transitivo \forma{pensar}; sin antecedente.
5. [\etq{Ocomp} Dieron un premio a [\etq{OR/CT} \textit{quienes}
participaron]]: \forma{quienes}, sujeto de la OR; sin antecedente. 6.
[\etq{Ocomp} El coche [\etq{OR/AS} \textit{que} está estacionado allí] es el
mío]: \forma{que}, sujeto de la OR; antecedente, \forma{coche}. 7. [\etq{Ocomp}
El libro [\etq{OR/AS} \textit{que} me prestaste] es muy interesante]:
\forma{que}, complemento objeto directo del verbo transitivo \forma{prestar};
antecedente, \forma{libro}. 8. [\etq{Ocomp} [\etq{OR/S} \textit{El que} llegó
tarde] no entró]: \forma{el que}, sujeto de la OR; sin antecedente. 9.
[\etq{Ocomp} Esa es la casa [\etq{OR/AS} \textit{cuya} puerta estaba pintada de
verde]]: \forma{cuya}, especificador de \forma{puerta}; antecedente,
\forma{casa}. 10. [\etq{Ocomp} Esta es la lapicera [\etq{OR/AS} con \textit{la
cual} escribí la carta]]: \forma{la cual}, complemento de la preposición
\forma{con} (valor instrumental); antecedente, \forma{lapicera}. 11.
[\etq{Ocomp} Hizo la preparación [\etq{OR/AC} \textit{como} le enseñaron]]:
\forma{como}, adjunto circunstancial de modo o manera del verbo
\forma{enseñar}; sin antecedente. 12. [\etq{Ocomp} La docente [\etq{OR/AS} con
\textit{quien} estudié inglés] era británica]: \forma{quien}, complemento de la
preposición \forma{con} (valor sociativo); antecedente, \forma{docente}. 13.
[\etq{Ocomp} La película [\etq{OR/AS} por \textit{la cual} vinimos al festival]
fue excelente]: \forma{la cual}, complemento o término de la preposición
\forma{por} (valor causal); antecedente, \forma{película}. 14. [\etq{Ocomp} Lo
voy a comprar [\etq{OR/AC} \textit{cuando} esté en oferta]]: \forma{cuando},
adjunto circunstancial de tiempo del verbo \forma{estar}; sin antecedente. 15.
[\etq{Ocomp} Los estudiantes [\etq{OR/AS} \textit{que} estudiaron] aprobaron el
examen]: \forma{que}, sujeto de la OR; antecedente, \forma{estudiantes}. 16.
[\etq{Ocomp} Los estudiantes, [\etq{OR/AX} \textit{que} estudiaron], aprobaron el
examen]: \forma{que}, sujeto de la OR; antecedente, \forma{estudiantes}. 17.
[\etq{Ocomp} Necesito [\etq{OR/COD} \textit{lo que} preparaste para el relleno
de la tarta]]: \forma{lo que}, complemento objeto directo del verbo transitivo
\forma{preparar}; sin antecedente. 18. [\etq{Ocomp} No me importa [\etq{OR/S}
\textit{lo que} digan]]: \forma{lo que}, complemento objeto directo del verbo
transitivo \forma{decir}; sin antecedente. 19. [\etq{Ocomp} [\etq{OR/S}
\textit{Quien} persevera], vence]: \forma{quien}, sujeto de la OR; sin
antecedente. 20. [\etq{Ocomp} [\etq{OR/S} \textit{Quienes} estudian], aprenden]:
\forma{quienes}, sujeto de la OR; sin antecedente.
\end{solucion}

\begin{solucion}{\ref{ej:13-paratextos}}
2. La OC encabezada por la conjunción \forma{que} (\forma{que no iba a volver a
hablar}) es el complemento objeto directo (COD) del verbo transitivo
\forma{pensar}. La OC encabezada por la conjunción \forma{que} (\forma{que está
feliz\ldots}) es el COD del verbo transitivo \forma{contar}. La OR encabezada
por el pronombre relativo \forma{que} (\forma{que dará en el escenario más
importante del país}) es un adjunto especificativo (AS) del antecedente nominal
\forma{concierto}. La OC (sin subordinante) con el pronombre interrogativo
\forma{cómo} (\forma{cómo fueron los momentos más duros de su enfermedad\ldots})
es el COD del verbo transitivo \forma{recordar}. La OR encabezada por el
pronombre relativo \forma{que} (\forma{que lo tuvo en coma dos meses}) es un
adjunto explicativo (AX) del antecedente nominal \forma{enfermedad}. La OS
encabezada por la conjunción \forma{que} (\forma{que es más lindo}) es un adjunto
circunstancial (AC) de causa (\sig{mejor me quedo y vivo, porque es más lindo}).
La oración completa entrecomillada en simples es una mención en estilo directo,
COD del verbo transitivo \forma{pensar}.

3. La OC de infinitivo (sin subordinante, \forma{tapar baches en 20 minutos}) es
el complemento o término (CT) de la preposición \forma{para}. La OR encabezada
por el pronombre relativo \forma{que} (\forma{que podían hacer las viejas
cuadrillas}) es un adjunto especificativo (AS) del antecedente elidido
\forma{aperturas} (\sig{la mitad de aperturas}). La OC encabezada por la
conjunción \forma{que} (\forma{que íbamos a tapar cada bache en no más de 15
días}) es el complemento o término (CT) de la preposición
\forma{a}.\pendiente{El original atribuye a \textit{a} «valor final» en
\textit{comprometerse a}}

4. La OR encabezada por el pronombre relativo \forma{que} (\forma{que
revolucionó al mundo del deporte}) es un adjunto especificativo (AS) del
antecedente nominal \forma{carta}. La OR encabezada por el pronombre relativo
\forma{que} (\forma{que así ganó la Supercopa de Francia}) es un adjunto
explicativo (AX) del antecedente nominal \forma{PSG}. La OS de gerundio (sin
subordinante, \forma{posteando un texto contra\ldots}) es un adjunto
circunstancial (AC) de modo o manera del verbo \forma{sorprender}
(alternativamente, la construcción acepta la interpretación temporal, \sig{al
postear un texto contra\ldots}). La OR semilibre encabezada por el pronombre
relativo \forma{los que} (\forma{los que critican a los futbolistas}) es el
complemento o término (CT) de la preposición \forma{contra}. La OC encabezada
por la conjunción \forma{que} (\forma{que el fútbol no tiene nada que ver con la
vida}) es el COD del verbo transitivo \forma{decir}. (Interpretamos que la
locución verbal \forma{tener que ver} se encuentra cuantificada en su grado,
\forma{tener \{nada / algo / bastante / mucho\} que ver}, con mantenimiento de
su régimen preposicional). La OC (sin subordinante) con el pronombre
interrogativo \forma{cuánto} (\forma{cuánto saben de la vida}) es el COD del
verbo transitivo \forma{saber}.

5. La OC encabezada por la conjunción \forma{si} (\forma{si un tipo de mi edad
se deja influenciar por alguien\ldots}) es el sujeto (S) de la oración.
(Interpretamos como perífrasis verbal \forma{dejarse influenciar}, con valor
modal radical, para indicar permiso). La OR encabezada por el pronombre relativo
\forma{que} (\forma{que está defendiendo sus cosas}) es un adjunto
especificativo (AS) del antecedente pronominal indefinido \forma{alguien}. La OR
encabezada por el adverbio relativo \forma{como} (\forma{como no lo hice
nunca}), cuyo antecedente es toda la oración anterior, constituye un adjunto
explicativo (AX) de esta (\sig{no voy a mandar al frente a nadie, de la misma
manera en que no lo hice nunca}). La OS encabezada por la conjunción subordinante
\forma{si} (\forma{si no me devuelven mi plata}) es un adjunto de nivel
oracional, un modificador de núcleo oracional (MNO) que introduce un estado de
cosas en el mundo que da las condiciones de verdad para el período hipotético
condicional (en términos ilocutivos, toda la oración representa una amenaza).

6. La OC de infinitivo (sin subordinante, \forma{volver a jugar otra final}) es
el predicativo subjetivo obligatorio (PSO) del verbo copulativo \forma{ser}. La
OC encabezada por la conjunción subordinante \forma{que} (\forma{que lo
merecemos}) es el complemento o término (CT) de la preposición \forma{por} (valor
causal).

7. La OR encabezada por el pronombre relativo \forma{el que} (\forma{en el que
protagoniza escenas jugadas}) es un adjunto especificativo (AS) del antecedente
nominal \forma{rol}. La OC encabezada por la conjunción \forma{que} (\forma{que
ese rol la puede reposicionar, tras\ldots}) es el COD del verbo \forma{sentir}.
La OC de infinitivo (sin subordinante, \forma{perder espacio para\ldots}) es
complemento o término (CT) de la preposición \forma{tras} (valor temporal). La
OC de infinitivo (sin subordinante, \forma{resguardar su intimidad}) es
complemento o término (CT) de la preposición \forma{para} (valor final). La OR
encabezada por el relativo \forma{que} (\forma{que cotiza alto en este oficio})
es un adjunto especificativo (AS) del antecedente nominal \forma{moneda}.

8. La OC encabezada por la conjunción \forma{que} (\forma{que si le regresan el
dinero\ldots}) es el COD del verbo \forma{asegurar}. La OS encabezada por la
conjunción subordinante \forma{si} (\forma{si le regresan el dinero\ldots}) es
un adjunto de nivel oracional, un modificador de núcleo oracional (MNO) que
introduce un estado de cosas en el mundo que da las condiciones de verdad para
el período hipotético condicional (en términos ilocutivos, toda la oración
representa una promesa). La OR encabezada por el pronombre relativo \forma{que}
(\forma{que le falta}) es un adjunto especificativo (AS) del antecedente nominal
\forma{dinero}. La OR libre encabezada por el pronombre relativo \forma{quien}
(\forma{quien ayer asistió a la primera audiencia}) es una aposición explicativa
(APX) del nombre propio, correferente y con el que puede cambiar posiciones, lo
que no sería esperable de una OR con antecedente nominal, que debe precederla.

9. La OC de infinitivo (sin subordinante, \forma{sumar un feriado en junio}) es
el complemento o término (CT) de la preposición \forma{de}. La OC de infinitivo
(sin subordinante, \forma{recordar la muerte del General Güemes\ldots}) es el
complemento de la preposición \forma{para}. La OS encabezada por la conjunción
subordinante \forma{si} (\forma{si obtiene la sanción definitiva del Congreso})
es un adjunto de nivel oracional, un modificador de núcleo oracional (MNO) que
introduce un estado de cosas en el mundo que da las condiciones de verdad para
el período hipotético condicional.

10. La OR encabezada por el pronombre relativo \forma{que} (\forma{que halló la
Policía}) es un adjunto especificativo (AS) del antecedente nominal
\forma{elementos}.

\pendiente{El original reproduce además cada paratexto con todos los corchetes;
se omite junto con los textos}
\end{solucion}
